\chapter{Complex spaces and modifications}
\label{chap:a01}

Analytic sets enter pluripotential theory as the pole sets of potentials with analytic singularities, as the zero sets of the ideals that describe these poles, and as the exceptional divisors of modifications. This appendix collects the facts about analytic sets, reduced complex spaces and modifications that are used in the book. Everything is proved, except for two existence theorems from the theory of resolution of singularities, which are stated in \cref{sec:analytic-blowups-resolution}.

The local theory rests on the Weierstrass division theorem. It shows that the ring of germs of holomorphic functions is Noetherian and, through Oka's theorem, that the sheaf of holomorphic functions is coherent, so that statements about one stalk extend to a neighborhood. The local parametrization theorem (\cref{thm:analytic-local-parametrization}) presents an irreducible germ of an analytic set as a finite branched covering of a polydisc. From it we deduce the analytic Nullstellensatz, the basic properties of dimension, Cartan's theorem on the coherence of the ideal sheaf of an analytic set, and the decomposition of an analytic set into irreducible components. One phenomenon requires care: an irreducible analytic set can have several local branches at one point. The normalization (\cref{thm:analytic-normalization}) separates these branches; its holomorphic functions are the weakly holomorphic functions, that is, the functions that are holomorphic on the regular locus and locally bounded near the singular locus. The cusp shows that a space with only one branch at each point need not be normal (\cref{ex:cusp-node-normalization}).

For modifications, explicit charts describe blowups along submanifolds and the transformation of ideals under them (\cref{prop:analytic-blowup-universal,prop:smooth-center-blowup}). Remmert's theorem, proved together with the Remmert--Stein extension theorem, shows that images of analytic sets under proper holomorphic maps, in particular the images of exceptional divisors, are analytic (\cref{thm:proper-analytic-image}). We then prove a relative principalization theorem, which controls an ideal and its restriction to a submanifold simultaneously (\cref{thm:relative-principalization}), show that proper bimeromorphic maps onto spaces with one local branch at each point have connected fibers (\cref{thm:unibranch-connected-fibers}), and prove that every prime divisor over a compact complex space appears on a model obtained by finitely many blowups with smooth centers (\cref{thm:divisor-smooth-blowup-realization}).

The last section defines plurisubharmonic functions on reduced complex spaces by local extension to an ambient space, proves that they are constant on compact connected spaces, and shows that they are locally bounded above near analytic subsets of codimension at least two. Holomorphic curves through singular points, constructed in \cref{sec:analytic-curves-finite-maps}, reduce these statements to the maximum principle in one variable. Together, these results supply the facts about complex spaces and modifications used in the later chapters of Part~II and in \cite{XiaSampleBook}. Historical remarks and references for this appendix are collected in \cref{sec:history-analytic-spaces}.

\section{Division and finite analytic modules}
\label{sec:analytic-division}

Let $n\geq0$. We write $\mathcal O_n$ for the ring of germs at $0\in\mathbb C^n$ of holomorphic functions; two functions defined near $0$ have the same germ if they agree on some neighborhood of $0$, and identities between germs are understood in this sense. By the iterated Cauchy formula of \cref{sec:complex-exterior-calculus}, germs are the same as convergent power series, so $\mathcal O_n=\mathbb C\{z_1,\ldots,z_n\}$, and $\mathcal O_0=\mathbb C$. By the identity theorem, $\mathcal O_n$ is an integral domain. A germ that does not vanish at the origin has a holomorphic reciprocal near the origin; hence $\mathcal O_n$ is a local ring whose maximal ideal consists of the germs vanishing at the origin. The \emph{order} of a nonzero germ is the degree of the first nonzero homogeneous term of its Taylor expansion; the order of a product is the sum of the orders, since the first homogeneous terms multiply.

Write points of $\mathbb C^n$, $n\geq1$, as $(z',w)$ with $z'\in\mathbb C^{n-1}$. A \emph{Weierstrass polynomial} of degree $m\geq0$ in $w$ is a polynomial
\[
 P(z',w)=w^m+a_1(z')w^{m-1}+\cdots+a_m(z'),\qquad a_j\in\mathcal O_{n-1},\quad a_j(0)=0.
\]
We shall often use the following bound for the roots $t$ of a monic polynomial with complex coefficients:
\begin{equation}\label{eq:monic-root-bound}
 t^m+d_1t^{m-1}+\cdots+d_m=0\quad\Longrightarrow\quad|t|\le2\max_{1\leq j\leq m}|d_j|^{1/j}.
\end{equation}
Indeed, otherwise $|d_j|<(|t|/2)^j$ for all $j$, and dividing the equation by $t^m$ gives $1\le\sum_{j=1}^m|d_j/t^j|<\sum_{j=1}^m2^{-j}<1$. In particular, the roots of $P(z',\cdot)$ tend to $0$ as $z'\to0$.

A germ $f\in\mathcal O_n$ is \emph{distinguished} in $w$ if $f(0,w)$ is not identically zero; if $f(0,w)$ vanishes to order $m\geq1$ at $w=0$, the Weierstrass preparation theorem (\cref{thm:local-weierstrass-preparation}) writes $f$ as a unit times a Weierstrass polynomial of degree $m$, and if $f(0)\neq0$, then $f$ is itself a unit, which is the case $m=0$. Every nonzero germ becomes distinguished in the last coordinate after a linear change of coordinates. Indeed, let $P_s$ be the first nonzero homogeneous term of $f$ and choose $v$ with $P_s(v)\neq0$; in linear coordinates whose last basis vector is $v$, we have $f(0,w)=P_s(v)w^s+O(w^{s+1})$, which vanishes to order exactly $s$, the order of $f$. A polynomial vanishing on a nonempty open set is zero, so the admissible vectors $v$ form a dense open set. Hence the new coordinates can be chosen arbitrarily close to given ones, and finitely many nonzero germs become distinguished simultaneously if $v$ avoids the zero set of the product of their first homogeneous terms.

\begin{proposition}[Weierstrass division]\label{prop:weierstrass-division}
Let $n\geq1$ and $m\geq1$, and let $P$ be a Weierstrass polynomial of degree $m$ in $w$. Then every germ $f\in\mathcal O_n$ can be written uniquely as
\[
 f=Pq+R,\qquad q\in\mathcal O_n,\quad
 R\in\mathcal O_{n-1}[w],\quad\deg_w R<m.
\]
Moreover, let $\Delta'\subset\mathbb C^{n-1}$ be a polydisc and $\Delta=\{|w|<r\}$ a disc such that the coefficients of $P$ are holomorphic near $\overline{\Delta'}$ and all roots of $P(z',\cdot)$, $z'\in\Delta'$, lie in $\{|w|\leq\rho\}$ for some $\rho<r$; such polydiscs $\Delta'\times\Delta$ exist with arbitrarily small radii. Then division is a bounded linear operation on bounded holomorphic functions on $\Delta'\times\Delta$: every such $f$ can be written uniquely as $f=Pq+R$, where $q$ is holomorphic on $\Delta'\times\Delta$ and $R$ is a polynomial in $w$ of degree less than $m$ with holomorphic coefficients on $\Delta'$, and there is a constant $C$, independent of $f$, with
\[
 \sup_{\Delta'\times\Delta}|q|+\sup_{\Delta'\times\Delta}|R|\le C\sup_{\Delta'\times\Delta}|f|.
\]
\end{proposition}
\begin{proof}
Polydiscs as in the statement exist with arbitrarily small radii, because the roots of $P(z',\cdot)$ tend to $0$ as $z'\to0$, by \eqref{eq:monic-root-bound}.

Let $\Delta'\times\Delta$ be such a polydisc and $f$ a bounded holomorphic function on it. For $\rho<s<r$ and $|w|<s$ define
\[
 q(z',w)=\frac1{2\pi\mathrm i}\int_{|\zeta|=s}
       \frac{f(z',\zeta)}{P(z',\zeta)(\zeta-w)}\,\mathrm d\zeta,
 \qquad
 R(z',w)=\frac1{2\pi\mathrm i}\int_{|\zeta|=s}
 \frac{f(z',\zeta)}{P(z',\zeta)}
       \frac{P(z',\zeta)-P(z',w)}{\zeta-w}\,\mathrm d\zeta.
\]
For fixed $z'$ and $w$, both integrands are holomorphic in $\zeta$ on the annulus $\{\rho<|\zeta|<r\}$, except for the pole of the first one at $\zeta=w$; by Cauchy's theorem the integrals do not depend on $s$ as long as $|w|<s$. Hence $q$ and $R$ are defined on $\Delta'\times\Delta$, and they are holomorphic by differentiation under the integral sign. The difference quotient $(P(z',\zeta)-P(z',w))/(\zeta-w)$ is a polynomial in $w$ of degree $m-1$ whose coefficients are polynomials in $\zeta$ and in the $a_j(z')$; hence $R$ is a polynomial in $w$ of degree less than $m$ with holomorphic coefficients on $\Delta'$. Adding the two integrals, we obtain $Pq+R=(2\pi\mathrm i)^{-1}\int_{|\zeta|=s}f(z',\zeta)(\zeta-w)^{-1}\,\mathrm d\zeta=f$ by Cauchy's formula.

For the estimate fix $s_0=(\rho+r)/2$. For $|\zeta|\geq s_0$ we have $|P(z',\zeta)|=\prod_l|\zeta-\lambda_l(z')|\geq(s_0-\rho)^m$, where $\lambda_l(z')$ are the roots, and the coefficients of the difference quotient are bounded for $z'\in\Delta'$, $|\zeta|=s_0$ and $|w|<r$, since the $a_j$ are bounded on $\Delta'$. The integral over $|\zeta|=s_0$ therefore gives $|R|\leq C_1\sup|f|$ on $\Delta'\times\Delta$. For $s_0\leq|w|<r$ we have $|q|=|f-R|/|P|\leq(1+C_1)(s_0-\rho)^{-m}\sup|f|$, and the maximum principle for the holomorphic functions $q(z',\cdot)$ on the discs $\{|w|\leq s\}$, $s_0\leq s<r$, extends this bound to all of $\Delta'\times\Delta$.

For uniqueness, suppose that $Pq+R=0$ on a polydisc as above, with $R$ a polynomial in $w$ of degree less than $m$. For each $z'$, the polynomial $R(z',\cdot)=-P(z',\cdot)q(z',\cdot)$ vanishes at the $m$ roots of $P(z',\cdot)$, counted with multiplicity, all of which lie in $\Delta$. Its degree is less than $m$, so $R(z',\cdot)=0$. Then $P(z',\cdot)q(z',\cdot)=0$, and since $P(z',\cdot)$ has only finitely many zeros, $q(z',\cdot)=0$ by continuity. Finally, every germ is represented by a bounded holomorphic function on some polydisc as above, and an identity between germs holds on some such polydisc; so the assertions for germs follow.
\end{proof}
\index{Weierstrass division}

An \emph{analytic subset} of a complex manifold $M$ is a closed subset $A\subset M$ such that every point of $M$ has a neighborhood $U$ and finitely many holomorphic functions $g_1,\ldots,g_k$ on $U$ with $A\cap U=\{g_1=\cdots=g_k=0\}$. It is \emph{proper} if $A\neq M$.
\index{analytic subset}

\begin{lemma}\label{lem:bounded-holomorphic-removal}
Let $A$ be a proper analytic subset of a connected complex manifold $M$. Then $A$ is nowhere dense in $M$, and $M\setminus A$ is connected. Every holomorphic function on $M\setminus A$ which is locally bounded near $A$ extends uniquely to a holomorphic function on $M$. Moreover, the ring $\mathcal O_n$ is integrally closed in its field of fractions.
\end{lemma}
\begin{proof}
The interior of $A$ is open; we show that it is also closed. Let $a\in M$ be a limit of interior points of $A$, and choose a connected coordinate neighborhood $U$ of $a$ and functions $g_1,\ldots,g_k$ with $A\cap U=\{g_1=\cdots=g_k=0\}$. Each $g_i$ vanishes on a nonempty open subset of $U$, hence on $U$ by the identity theorem; so $U\subset A$ and $a$ is an interior point. Since $M$ is connected and $A\neq M$, the interior of $A$ is empty. Consequently, for every $a\in A$ and $U$ as above, some $g_i$ is not identically zero on $U$, and $A\cap U$ is contained in the zero set of a holomorphic function $g\not\equiv0$ on $U$.

Let $f$ be holomorphic on $M\setminus A$ and locally bounded near $A$. If $\dim M=0$, then $A=\varnothing$ and there is nothing to prove. Fix $a\in A$ and $g$ as above, and choose linear coordinates $(z',w)$ centered at $a$ in which $g$ is distinguished in $w$. By the Weierstrass preparation theorem and \eqref{eq:monic-root-bound}, there is a polydisc $\Delta'\times\Delta$ centered at $a$, with $\Delta=\{|w|<r\}$, on which $g$ is a unit times a Weierstrass polynomial $W$ and $f$ is bounded off $A$, such that all roots of $W(z',\cdot)$, $z'\in\Delta'$, have modulus at most some $\rho<r$. For fixed $z'\in\Delta'$, the function $f(z',\cdot)$ is holomorphic on $\Delta$ outside the finitely many points of $A$ in the fiber, which are roots of $W(z',\cdot)$, and it is bounded near them; by the removable singularity theorem in one variable it extends holomorphically to $\Delta$. For $\rho<s<r$ the function
\[
 F(z',w)=\frac1{2\pi\mathrm i}\int_{|\zeta|=s}\frac{f(z',\zeta)}{\zeta-w}\,\mathrm d\zeta,\qquad z'\in\Delta',\ |w|<s,
\]
is holomorphic, by differentiation under the integral sign, since $f$ is holomorphic near $\Delta'\times\{|\zeta|=s\}$, a set disjoint from $A$. By Cauchy's formula for the extended functions $f(z',\cdot)$, it agrees with $f$ outside $A$. Letting $s\uparrow r$, we obtain a holomorphic extension of $f$ to $\Delta'\times\Delta$. Since $A$ is nowhere dense, two holomorphic extensions on overlapping open sets agree, by continuity; so the local extensions glue to a holomorphic function on $M$, and this extension is unique.

If $M\setminus A$ were disconnected, the function equal to $0$ on one nonempty union of its connected components and to $1$ on the others would be bounded and holomorphic on $M\setminus A$. Its holomorphic extension takes values in $\{0,1\}$, by continuity and density of $M\setminus A$; since $M$ is connected, it is constant, a contradiction.

Finally, let $b/c$, with $b,c\in\mathcal O_n$ and $c\neq0$, satisfy a monic equation with coefficients in $\mathcal O_n$. On a small connected neighborhood $U$ of the origin, these coefficients are bounded, and $b/c$ is a holomorphic function on $U\setminus\{c=0\}$ satisfying the equation there; by \eqref{eq:monic-root-bound} it is bounded. By the first part, applied to the proper analytic subset $\{c=0\}$ of $U$, it extends to a holomorphic function $h$ on $U$, and $ch=b$. Hence $b/c=h\in\mathcal O_n$, which proves that $\mathcal O_n$ is integrally closed.
\end{proof}
\index{holomorphic function!bounded removal}\index{integrally closed ring}

\Needspace{9\baselineskip}
\begin{proposition}\label{prop:analytic-noetherian}
Every ideal of $\mathcal O_n$ is finitely generated, and every submodule of $\mathcal O_n^r$, $r\geq1$, is finitely generated. Consequently, every increasing sequence of ideals of $\mathcal O_n$, or of submodules of a finitely generated $\mathcal O_n$-module, is eventually constant.
\end{proposition}
\begin{proof}
We first recall three algebraic facts about a commutative ring $\Lambda$ all of whose ideals are finitely generated. Every submodule $N$ of $\Lambda^r$ is finitely generated, by induction on $r$: the projection of $N$ to the last coordinate is an ideal; choose elements of $N$ whose last coordinates generate it, and add generators of $N\cap(\Lambda^{r-1}\times\{0\})$, which exist by the induction hypothesis. A finitely generated module is a quotient of some $\Lambda^r$, and the inverse image of a submodule is a submodule of $\Lambda^r$; hence submodules of finitely generated modules are finitely generated. Finally, the union of an increasing sequence of submodules is a submodule; its finitely many generators already lie in one member of the sequence, which is therefore stationary from then on.

It remains to prove that the ideals of $\mathcal O_n$ are finitely generated, which we do by induction on $n$. The ring $\mathcal O_0=\mathbb C$ is a field. Let $n\geq1$ and let $I\subset\mathcal O_n$ be an ideal. If $I=0$ or $I$ contains a unit, it is generated by one element. Otherwise choose a nonzero germ in $I$, and coordinates in which it is distinguished in $w$. By the preparation theorem it is a unit times a Weierstrass polynomial $P$ of some degree $m\geq1$, and $P\in I$. Divide every $f\in I$ by $P$ using \cref{prop:weierstrass-division}; the remainder $R_f=f-Pq$ again belongs to $I$. The remainders of degree less than $m$ that lie in $I$ form an $\mathcal O_{n-1}$-submodule of $\mathcal O_{n-1}[w]$, which we identify, through their coefficients, with a submodule of $\mathcal O_{n-1}^m$. By the induction hypothesis and the first paragraph, it is generated by finitely many remainders $R_1,\ldots,R_k$. Then $P,R_1,\ldots,R_k$ generate $I$.
\end{proof}
\index{Noetherian ring!analytic germs}

We use sheaves in the following elementary form. A sheaf $\mathcal F$ of modules on a topological space assigns to each open set $U$ a module $\mathcal F(U)$ of sections, with restriction maps, such that sections that agree locally are equal and compatible local sections glue. The stalk $\mathcal F_x$ consists of the germs at $x$ of sections defined near $x$. Kernels and images of morphisms of sheaves are defined stalkwise, and a sequence of morphisms is exact if it is exact on every stalk; in particular, a section of the image of a morphism is, near each point, the image of a section of the source. Let $\mathcal A$ be a sheaf of rings and $\mathcal F$ a sheaf of $\mathcal A$-modules. We say that $\mathcal F$ is \emph{locally finitely generated} if every point has a neighborhood $U$ and finitely many sections $s_1,\ldots,s_p\in\mathcal F(U)$ whose germs generate $\mathcal F_y$ over $\mathcal A_y$ for every $y\in U$. This requires one family of generators on a whole neighborhood, not merely finite generation of each stalk. The \emph{sheaf of relations} of sections $s_1,\ldots,s_p\in\mathcal F(U)$ is the kernel of the morphism $\mathcal A^p|_U\to\mathcal F|_U$, $(a_i)\mapsto\sum_ia_is_i$. The sheaf $\mathcal F$ is \emph{coherent} if it is locally finitely generated and the sheaf of relations of every finite family of sections over every open set is locally finitely generated.
\index{sheaf}\index{coherent sheaf}

\begin{lemma}\label{lem:coherent-sheaf-algebra}
Let $\mathcal A$ be a sheaf of rings. Kernels, images and cokernels of morphisms between coherent $\mathcal A$-modules are coherent. If two of the three sheaves in an exact sequence $0\to\mathcal F\to\mathcal S\to\mathcal G\to0$ of $\mathcal A$-modules are coherent, so is the third. Finite direct sums of coherent sheaves are coherent. If $\mathcal A$ is coherent as a module over itself, every locally finitely generated $\mathcal A$-submodule of $\mathcal A^r$ is coherent.
\end{lemma}
\begin{proof}
A locally finitely generated submodule of a coherent sheaf is coherent, because its sections have the same relations when they are regarded as sections of the larger sheaf. The image of a morphism from a coherent sheaf is locally generated by the images of local generators of the source, so the image of a morphism between coherent sheaves is coherent. For the kernel of a morphism $\phi\colon\mathcal F\to\mathcal G$ between coherent sheaves, choose local generators $F_1,\ldots,F_p$ of $\mathcal F$ and local generators $R_j=(R_{j1},\ldots,R_{jp})$ of the relations of $\phi(F_1),\ldots,\phi(F_p)$. If $\sum_ia_iF_i$ lies in the kernel, then $(a_i)$ is a relation of the $\phi(F_i)$, hence a combination of the $R_j$; so the sections $\sum_iR_{ji}F_i$ generate the kernel, which is therefore coherent.

Let $\mathcal F\subset\mathcal S$ be coherent and consider $\mathcal G=\mathcal S/\mathcal F$; this covers cokernels, because images are coherent. Local generators of $\mathcal S$ map to local generators of $\mathcal G$. Given sections $G_1,\ldots,G_p$ of $\mathcal G$, lift them locally to sections $S_1,\ldots,S_p$ of $\mathcal S$, and choose local generators $F_1,\ldots,F_k$ of $\mathcal F$. A tuple $(a_i)$ is a relation of the $G_i$ exactly when $\sum_ia_iS_i\in\mathcal F$, that is, when $(a,b)$ is a relation of $(S_1,\ldots,S_p,F_1,\ldots,F_k)$ in $\mathcal S$ for some $b$. Hence the projections to the first $p$ coordinates of local generators of the latter relations generate the relations of the $G_i$, and $\mathcal G$ is coherent.

Now let $\mathcal F$ and $\mathcal G$ be coherent. Local lifts of generators of $\mathcal G$, together with local generators of $\mathcal F$, generate $\mathcal S$. Given sections $S_1,\ldots,S_p$ of $\mathcal S$, let $R_1,\ldots,R_l$ generate the relations of their images in $\mathcal G$. The combinations $H_j=\sum_iR_{ji}S_i$ lie in $\mathcal F$; let $Q_1,\ldots,Q_k$ generate their relations. Every relation $(a_i)$ of the $S_i$ is first a relation of their images, so $a=\sum_jc_jR_j$, and then $\sum_jc_jH_j=\sum_ia_iS_i=0$, so $c$ is a combination of the $Q_k$. Hence the tuples $\sum_jQ_{kj}R_j$, which are relations of the $S_i$, generate all of them, and $\mathcal S$ is coherent. The remaining case of the exact sequence, with $\mathcal S$ and $\mathcal G$ coherent, is the kernel case. A finite direct sum is an extension $0\to\mathcal F\to\mathcal F\oplus\mathcal G\to\mathcal G\to0$. Finally, if $\mathcal A$ is coherent, so is $\mathcal A^r$, and the last assertion follows from the first sentence of the proof.
\end{proof}

\begin{theorem}[Oka's coherence theorem]\label{thm:oka-coherence}
The sheaf $\mathcal O_M$ of holomorphic functions on a complex manifold $M$ is coherent. In particular, the kernel of the morphism $\mathcal O_U^q\to\mathcal O_U^p$ defined by a $p\times q$ matrix of holomorphic functions on an open set $U\subset M$ is locally finitely generated.
\end{theorem}
\begin{proof}
The sheaf $\mathcal O_M$ is generated by the constant $1$, and the statement is local, so we must show that the relations of finitely many holomorphic functions $F_1,\ldots,F_q$ near a point of $\mathbb C^n$, which we take to be the origin, are locally finitely generated. We argue by induction on $n$. For $n=0$ this is finite-dimensional linear algebra. Let $n\geq1$. If the germ of some $F_j$ at $0$ is zero, then $F_j$ vanishes on a connected neighborhood, and the relations of the $F_i$ are those of the remaining functions, with an arbitrary $j$th coordinate; so we may assume that no germ $F_j$ is zero (if all are zero, every tuple is a relation). Choose coordinates $(z',w)$ in which all $F_j$ are distinguished in $w$. By the preparation theorem, on a small polydisc each $F_j$ is a holomorphic unit $u_j$ times a Weierstrass polynomial; since $(g_j)$ is a relation of the $F_j$ exactly when $(u_jg_j)$ is a relation of the Weierstrass polynomials, we may assume that each $F_j$ is a Weierstrass polynomial in $w$ whose coefficients are holomorphic on a polydisc $\Delta'\subset\mathbb C^{n-1}$. Renumbering, let $F_q$ have the maximal degree $\mu$.

We claim that at every point $x=(x',w_0)$ of the polydisc, the stalk of relations is generated by relations whose coordinates are polynomials in $w$ of degree at most $\mu$ with coefficients holomorphic near $x'$. If $F_q(x)\neq0$, then for a relation $g=(g_1,\ldots,g_q)$ we have $g_q=-\sum_{j<q}g_jF_j/F_q$, so $g=\sum_{j<q}(g_j/F_q)(F_qe_j-F_je_q)$ is a combination of the elementary relations $F_qe_j-F_je_q$, which are polynomial of degree at most $\mu$. Suppose that $F_q(x)=0$, and let $\mu'\geq1$ be the order of the root $w_0$ of $F_q(x',\cdot)$. By the preparation theorem at $x$, $F_q$ is a unit near $x$ times a Weierstrass polynomial $P_x$ of degree $\mu'$ in $w-w_0$. Dividing the monic polynomial $F_q$ by $P_x$ as polynomials in $w$, and using the uniqueness in \cref{prop:weierstrass-division} at $x$, we see that this unit is a polynomial $Q_x$ in $w$ of degree $\mu-\mu'$, so that
\[
 F_q=P_xQ_x,\qquad Q_x(x)\neq0.
\]
Let $g$ be a relation at $x$. Divide each $g_j$, $j<q$, by $P_x$, and absorb $Q_x^{-1}$ in the quotient: $g_j=F_qt_j+r_j$ with $\deg_wr_j<\mu'$. Subtracting $\sum_{j<q}t_j(F_qe_j-F_je_q)$ from $g$, we obtain a relation $(r_1,\ldots,r_q)$ with $r_1,\ldots,r_{q-1}$ as above, and
\[
 P_xQ_xr_q=-\sum_{j<q}F_jr_j.
\]
The right-hand side is a polynomial in $w$; dividing it by $P_x$ as a polynomial and comparing with the uniqueness of division at $x$, we find that $Q_xr_q$ is a polynomial of degree less than $\mu$. Each $Q_xr_j$, $j<q$, is also a polynomial of degree at most $(\mu-\mu')+(\mu'-1)<\mu$. Thus $(Q_xr_1,\ldots,Q_xr_q)$ is a polynomial relation of degree less than $\mu$, and $(r_j)$ is its product with the unit $Q_x^{-1}$. This proves the claim.

A tuple of polynomials $g_j=\sum_{k=0}^{\mu}c_{jk}(z')w^k$ is a relation exactly when the $2\mu+1$ coefficients of the polynomial $\sum_jF_jg_j$ vanish. These coefficients are given by a holomorphic matrix acting on the coefficient vectors $(c_{jk})$, with entries formed from the coefficients of the $F_j$. By the induction hypothesis and \cref{lem:coherent-sheaf-algebra}, the kernel of this matrix is coherent, so it has finitely many generators on a smaller polydisc $\Delta'_1$, generating its stalk at every point of $\Delta'_1$. The corresponding polynomial tuples are relations of the $F_j$, and they generate, over the rings of germs on the base, all polynomial relations of degree at most $\mu$ at every point over $\Delta'_1$. By the claim, they generate all relations at every point of the polydisc over $\Delta'_1$. This proves that $\mathcal O_M$ is coherent. The last assertion follows from \cref{lem:coherent-sheaf-algebra}, since the kernel of a morphism between the coherent sheaves $\mathcal O_U^q$ and $\mathcal O_U^p$ is coherent.
\end{proof}
\index{Oka coherence theorem}

\Needspace{10\baselineskip}
\section{Analytic sets and their local branches}
\label{sec:analytic-local-branches}

Let $M$ be a complex manifold and $x\in M$. Two subsets of $M$ define the same \emph{germ} at $x$ if they agree on a neighborhood of $x$; an \emph{analytic germ} at $x$ is the germ $(A,x)$ of an analytic subset $A$ of a neighborhood of $x$. Holomorphic coordinates identify the ring $\mathcal O_{M,x}$ of germs of holomorphic functions at $x$ with $\mathcal O_n$. For an analytic germ $(A,x)$ let $I_{A,x}\subset\mathcal O_{M,x}$ be the ideal of germs vanishing on $(A,x)$. Conversely, for an ideal $I\subset\mathcal O_{M,x}$ let $V(I)$ be the germ at $x$ of the common zero set of finitely many generators of $I$, which exist by \cref{prop:analytic-noetherian}; it does not depend on the choice of generators. Both operations reverse inclusions, and $V(I_{A,x})=(A,x)$, because local defining functions of $A$ belong to $I_{A,x}$. Finite unions and finite intersections of analytic subsets are analytic: they are defined locally by the products of the defining functions and by the union of the families of defining functions, respectively. A nonempty analytic germ is \emph{irreducible} if it is not the union of two analytic subgerms different from it.

A point $a$ of an analytic subset $A\subset M$ is \emph{regular} if $A$ is a complex submanifold of $M$ near $a$; the dimension of this submanifold is the dimension of $A$ at the regular point $a$. The regular points form an open subset $A_{\mathrm{reg}}$ of $A$, the \emph{regular locus}, and $A_{\mathrm{sing}}=A\setminus A_{\mathrm{reg}}$ is the \emph{singular locus}.
\index{analytic subset!regular locus}

\begin{proposition}\label{prop:analytic-irreducible-germs}
Every analytic germ is the union of finitely many irreducible analytic germs, and this decomposition is unique up to order if none of its members is contained in another. A nonempty analytic germ is irreducible if and only if its ideal is prime.
The members of the decomposition are called the \emph{irreducible branches}, or \emph{local branches}, of the germ.
\end{proposition}
\begin{proof}
A decreasing sequence of analytic germs is eventually constant, because the corresponding ideals increase and are eventually constant by \cref{prop:analytic-noetherian}, and $V(I_{A,x})=(A,x)$. Split a reducible germ into two analytic subgerms different from it, split each reducible piece again, and so on. If this process did not terminate, the resulting tree, in which every vertex has at most two children, would be infinite, and choosing at each level a child with infinitely many descendants would produce a strictly decreasing infinite sequence of analytic germs, which is impossible. Hence a finite decomposition into irreducible germs exists; removing members contained in others, we may assume that no member is contained in another. Note that an irreducible germ contained in a finite union of analytic germs is contained in one of them, since it is the union of its intersections with them. Now let $\bigcup_iA_i=\bigcup_jB_j$ be two such decompositions. Each $A_i$ is contained in some $B_j$, which in turn is contained in some $A_k$; then $A_i\subset A_k$ forces $k=i$ and $A_i=B_j$. By symmetry the two decompositions agree.

Let $(A,x)$ be irreducible. Its ideal is proper, since $x\in A$. If $fg$ vanishes on $(A,x)$, then $(A,x)$ is the union of its intersections with $\{f=0\}$ and $\{g=0\}$, so one of them is all of $(A,x)$, and $f$ or $g$ belongs to $I_{A,x}$; thus $I_{A,x}$ is prime. Conversely, suppose that $(A,x)$ is the union of two analytic subgerms $(B,x)$ and $(C,x)$ different from it. Since $V(I_{B,x})=(B,x)\neq(A,x)$, there is a germ $f\in I_{B,x}\setminus I_{A,x}$, and similarly $g\in I_{C,x}\setminus I_{A,x}$. Then $fg$ vanishes on $(A,x)$ while neither factor does, so $I_{A,x}$ is not prime.
\end{proof}
\index{analytic subset!irreducible germ}\index{analytic subset!local branches}

\begin{lemma}\label{lem:finite-integral-algebra}
Let $A\subset B$ be integral domains such that $A$ is integrally closed and $B$ is finitely generated as an $A$-module, and let $K=\operatorname{Frac}(A)$. Then $L=\operatorname{Frac}(B)$ is a finite extension of $K$. Suppose in addition that $A$ contains $\mathbb C$ and that $b_1,\ldots,b_r$ generate $B$ as an $A$-algebra. Then for all $(c_1,\ldots,c_r)\in\mathbb C^r$ outside finitely many proper linear subspaces, the element $u=\sum_jc_jb_j$ generates $L$ over $K$, and its minimal polynomial $W$ over $K$ is monic with coefficients in $A$. Let $q=\deg W$, and let $\delta=\prod_{i<k}(\tau_i-\tau_k)^2$ be the discriminant of $W$, where $\tau_1,\ldots,\tau_q$ are the roots of $W$ in a splitting field. Then $\delta\in A\setminus\{0\}$, and every element $g\in L$ that is integral over $A$ satisfies
\[
 \delta g=\sum_{j=0}^{q-1}d_ju^j,\qquad d_j\in A.
\]
\end{lemma}
\begin{proof}
We first recall the determinant criterion for integrality. Let $C$ be a domain containing $A$, let $b\in C$, and suppose that multiplication by $b$ maps a submodule generated by $v_1,\ldots,v_N$, not all zero, into itself: $bv_i=\sum_ja_{ij}v_j$ with $a_{ij}\in A$. Multiplying the matrix $bI-(a_{ij})$ by its adjugate shows that the monic polynomial $\det(bI-(a_{ij}))$ in $b$ annihilates every $v_i$; since $C$ is a domain, this determinant is zero, and $b$ is integral over $A$. In particular every element of $B$ is integral over $A$. Conversely, if $x_1,\ldots,x_k\in C$ are integral over $A$, then $A[x_1,\ldots,x_k]$ is a finitely generated $A$-module, since powers of each $x_i$ can be reduced with its monic equation; by the criterion, sums and products of integral elements are integral. Likewise, an element integral over a ring generated by integral elements is integral over $A$: adjoin it together with the finitely many coefficients of its equation.

The ring $KB=\{b/a:b\in B,\ a\in A\setminus\{0\}\}$ is a domain and a finite-dimensional vector space over $K$. Multiplication by a nonzero element is an injective linear map of $KB$, hence surjective; so $KB$ is a field, and it equals $L$. Let $q=[L:K]$. In characteristic zero an irreducible polynomial $p$ has no repeated roots: its derivative is nonzero of smaller degree, so $p$ does not divide $p'$, and the Euclidean algorithm writes $1$ as a combination of $p$ and $p'$. Let $N$ be a splitting field over $K$ of the product of the minimal polynomials of $b_1,\ldots,b_r$. Extending embeddings one generator at a time, and using at each step that the minimal polynomial of $b_j$ over $K(b_1,\ldots,b_{j-1})$ divides that over $K$ and therefore has distinct roots in $N$, we obtain exactly $q$ distinct $K$-embeddings $\sigma_1,\ldots,\sigma_q$ of $L$ into $N$. Two distinct embeddings differ on some generator $b_j$, so for $i\neq k$ the set of $c\in\mathbb C^r$ with $\sum_jc_j(\sigma_i(b_j)-\sigma_k(b_j))=0$ is a proper linear subspace of $\mathbb C^r$. For $c$ outside these finitely many subspaces the $q$ values $\sigma_i(u)$ are distinct. They are roots of the minimal polynomial of $u$, whose degree $[K(u):K]$ is at most $q$; hence $K(u)=L$ and $W=\prod_i(T-\sigma_i(u))$. The coefficients of $W$ are elementary symmetric functions of the elements $\sigma_i(u)$, which are integral over $A$ because they satisfy the monic equation of $u$; so they are integral over $A$ and belong to $K$, hence to $A$. For the same reasons $\delta=\prod_{i<k}(\sigma_i(u)-\sigma_k(u))^2$ belongs to $A$, and it is nonzero.

Let $g\in L$ be integral over $A$, and write $g=\sum_{j<q}e_ju^j$ with $e_j\in K$. Applying the embeddings gives the linear system $\sigma_i(g)=\sum_je_j\sigma_i(u)^j$, $1\leq i\leq q$, whose Vandermonde determinant $D=\prod_{i<k}(\sigma_k(u)-\sigma_i(u))$ satisfies $D^2=\delta$. By Cramer's rule, $\delta e_j=D\cdot D_j$, where $D_j$ is the determinant obtained by replacing the $j$th column with the values $\sigma_i(g)$. Thus $\delta e_j$ is a polynomial with integer coefficients in the integral elements $\sigma_i(u)$ and $\sigma_i(g)$; it is integral over $A$ and belongs to $K$, so $\delta e_j\in A$. This proves the formula with $d_j=\delta e_j$.
\end{proof}
\index{integral extension!finite module}

\Needspace{24\baselineskip}
\begin{theorem}[Local parametrization]\label{thm:analytic-local-parametrization}
Let $(A,0)\subset(\mathbb C^n,0)$ be an irreducible analytic germ. There are linear coordinates $z=(z',z'')\in\mathbb C^d\times\mathbb C^{n-d}$, which can be chosen arbitrarily close to any given linear coordinates, an integer $q\geq1$, a representative $A$ of the germ and arbitrarily small polydiscs $\Delta'\times\Delta''$ with the following properties. The projection
\[
 \pi\colon A\cap(\Delta'\times\Delta'')\longrightarrow\Delta',\qquad
 \pi(z',z'')=z',
\]
is proper; its fiber over $0$ is $\{0\}$, and every fiber has at most $q$ points. There is a holomorphic function $\delta\not\equiv0$ on $\Delta'$ such that, over $\Delta'\setminus\{\delta=0\}$, the map $\pi$ is a connected $q$-sheeted covering whose sheets are graphs of holomorphic maps. Its total space $A_\delta=A\cap\pi^{-1}(\Delta'\setminus\{\delta=0\})$, the \emph{covering locus}, is a complex submanifold of dimension $d$, open and dense in $A\cap(\Delta'\times\Delta'')$.
\end{theorem}
\begin{proof}
We prove more generally the following statement for every prime ideal $J\subset\mathcal O_n$: the germ $V(J)$ has coordinates and polydiscs as in the theorem, and $I_{V(J),0}=J$. The theorem then follows by taking $J=I_{A,0}$, which is prime by \cref{prop:analytic-irreducible-germs} and satisfies $V(J)=(A,0)$. If $J=0$, then $V(J)$ is a neighborhood of $0$; we take $d=n$, $q=1$ and $\delta=1$, and $I_{V(0),0}=0$ by the identity theorem. Assume that $J\neq0$.

\emph{Step 1: triangular equations.} Since $J$ is prime, it is a proper ideal, so its elements vanish at $0$. Choose a nonzero element of $J$, of order $s_n\geq1$, and linear coordinates in which it is distinguished in $z_n$. By the preparation theorem it is a unit times a Weierstrass polynomial $P_n\in J$ of degree $s_n$ in $z_n$. Since a unit has order $0$, the polynomial $P_n$ also has order $s_n$, so its coefficient of $z_n^{s_n-j}$ has order at least $j$ in $(z_1,\ldots,z_{n-1})$. The ideal $J\cap\mathcal O_{n-1}$ of germs in the first $n-1$ coordinates is again prime. If it is nonzero, we repeat the construction with a linear change of $z_1,\ldots,z_{n-1}$ only, which preserves the properties of $P_n$; we continue until $J\cap\mathcal O_d=0$ for some $d$. This gives coordinates and Weierstrass polynomials
\[
 P_k=z_k^{s_k}+\sum_{j=1}^{s_k}a_{jk}(z_1,\ldots,z_{k-1})z_k^{s_k-j}\in J,
 \qquad a_{jk}=O(|(z_1,\ldots,z_{k-1})|^j),\qquad d<k\leq n.
\]
At each step only the zero set of a nonzero homogeneous polynomial is excluded in choosing the direction, so the coordinates can be chosen arbitrarily close to given ones. Fix generators of $J$, holomorphic on a neighborhood $U$ of $0$, and let $V(J)$ also denote their common zero set in $U$. Applying \eqref{eq:monic-root-bound} successively to $P_{d+1},\ldots,P_n$, we find a constant $C$ with
\[
 |z''|\leq C|z'|\qquad\text{for } z\in V(J) \text{ near } 0.
\]
Consequently, if $\Delta''$ is a small polydisc and $\Delta'$ a polydisc whose radius is small compared with that of $\Delta''$, the points of $V(J)\cap(\Delta'\times\Delta'')$ over a compact set $K\subset\Delta'$ lie in a compact subset of $K\times\Delta''$. Hence $\pi$ is proper, its fiber over $0$ is $\{0\}$, and its fibers are finite, since they lie in the sets of successive roots of $P_{d+1},\ldots,P_n$.

Successive division by $P_n,\ldots,P_{d+1}$ (\cref{prop:weierstrass-division}) shows that every germ is congruent modulo $J$ to a polynomial in $z_{d+1},\ldots,z_n$ of degree less than $s_k$ in $z_k$, with coefficients in $\mathcal O_d$. Since $J\cap\mathcal O_d=0$, the ring $\mathcal O_d$ injects into
\[
 B=\mathcal O_n/J,
\]
and $B$ is a finitely generated $\mathcal O_d$-module, generated by the classes of the monomials $\prod_{k>d}z_k^{\alpha_k}$ with $0\le\alpha_k<s_k$.

\emph{Step 2: a primitive element and division by the discriminant.} The ring $B$ is a domain, because $J$ is prime, and $\mathcal O_d$ is integrally closed by \cref{lem:bounded-holomorphic-removal}. We apply \cref{lem:finite-integral-algebra} with the generators $z_{d+1},\ldots,z_n$. Choose a primitive linear form $u=\sum_{k>d}c_kz_k$, let $W(z',T)\in\mathcal O_d[T]$ be its minimal polynomial, of degree $q$, and let $\delta\in\mathcal O_d\setminus\{0\}$ be its discriminant. The class of each $z_k$, $k>d$, is integral over $\mathcal O_d$, so its minimal polynomial $W_k(z',T)$ over $\operatorname{Frac}(\mathcal O_d)$ is monic of degree at most $q$ with coefficients in $\mathcal O_d$. These polynomials, and $W$, are Weierstrass polynomials in $T$. Indeed, by the preparation theorem in the variables $(z',T)$ at the origin, $W_k$ is a unit times a Weierstrass polynomial $V_k$ in $T$; dividing the monic polynomial $W_k$ by $V_k$ as polynomials and using the uniqueness in \cref{prop:weierstrass-division}, we see that this unit is a polynomial $Q_k$ in $T$. Substituting $T=z_k$, which vanishes at the origin, $Q_k(z',z_k)$ becomes a unit of $\mathcal O_n$, so $V_k(z',z_k)\in J$. By minimality $\deg V_k\geq\deg W_k$, hence $W_k=V_k$. The same argument applies to $W$ and $u$. By \cref{lem:finite-integral-algebra} there are polynomials $B_k(z',T)\in\mathcal O_d[T]$ of degree less than $q$ with
\[
 \delta z_k\equiv B_k(z',u)\pmod J,\qquad d<k\leq n.
\]
Let $G\subset J$ be the ideal generated by $W(z',u)$ and the germs $\delta z_k-B_k(z',u)$. We claim that, with $m=\max\{q,(n-d)(q-1)\}$, every germ $f\in\mathcal O_n$ satisfies
\begin{equation}\label{eq:analytic-discriminant-division}
 \delta^m f\equiv R_f(z',u)\pmod G,
 \qquad f\in J\Longrightarrow R_f=0,
\end{equation}
for a unique polynomial $R_f(z',T)\in\mathcal O_d[T]$ of degree less than $q$. To prove this, divide $f$ successively by the Weierstrass polynomials $W_n(z',z_n),\ldots,W_{d+1}(z',z_{d+1})$, whose coefficients depend only on $z'$. This writes $f$ as a polynomial $\sum_\alpha e_\alpha(z')z''^\alpha$ with $\alpha_k\leq q-1$ for all $k$, plus multiples of the $W_k(z',z_k)$. Multiplying by $\delta^m$ and substituting $\delta z_k\equiv B_k(z',u)$ modulo $G$, which is possible because $|\alpha|\leq(n-d)(q-1)\leq m$, we obtain a polynomial in $u$ with coefficients in $\mathcal O_d$. The discarded multiples of $W_k(z',z_k)$ disappear modulo $G$, because $\delta^qW_k(z',z_k)\in G$ and $m\geq q$. Indeed, substituting $\delta z_k\equiv B_k(z',u)$ in $\delta^qW_k(z',z_k)$ gives $H_k(z',u)$ modulo $G$ for some $H_k\in\mathcal O_d[T]$; since $H_k(z',u)\equiv\delta^qW_k(z',z_k)\equiv0$ modulo $J$, the polynomial $H_k$ vanishes at the class of $u$ in $\operatorname{Frac}(B)$, so it is divisible by the monic polynomial $W$ in $\mathcal O_d[T]$, and $H_k(z',u)\in G$. Finally, dividing the resulting polynomial in $u$ by $W$ gives $R_f$. If two polynomials of degree less than $q$ were both admissible, their difference would vanish at the class of $u$, whose minimal polynomial has degree $q$; this proves uniqueness. If $f\in J$, then $R_f(z',u)\in J$, so $R_f$ vanishes at the class of $u$ and $R_f=0$.

\emph{Step 3: the covering off the discriminant.} Choose polydiscs as in Step 1 so small that all the germs above, the generators of $J$ and the identities \eqref{eq:analytic-discriminant-division} for these generators are defined near the closure of $\Delta'\times\Delta''$. Let $z'\in\Delta'$ with $\delta(z')\neq0$. Then $W(z',\cdot)$ has $q$ distinct roots $\tau_1,\ldots,\tau_q$, and the points of $V(G)$ over $z'$ are exactly the $q$ points
\[
 W(z',T)=0,\qquad z_k=B_k(z',T)/\delta(z'),\qquad T=\tau_1,\ldots,\tau_q.
\]
Indeed, a point of $V(G)$ satisfies these equations with $T=u(z'')$. Conversely, the polynomial $\sum_kc_kB_k(z',T)-\delta(z')T$ vanishes at the class of $u$, so it is divisible by $W$ and vanishes at each $\tau_i$; hence the point defined by $\tau_i$ satisfies $u(z'')=\tau_i$ and lies in $V(G)$. These $q$ points are distinct, because their values of $u$ are distinct, and by the holomorphic implicit function theorem the simple roots $\tau_i$, and hence the points, depend holomorphically on $z'$ locally. Similarly, $\delta^{\deg W_k}W_k(z',B_k(z',T)/\delta)$ is a polynomial in $T$ vanishing at the class of $u$, so the coordinate $z_k=B_k(z',\tau_i)/\delta(z')$ is a root of the Weierstrass polynomial $W_k(z',\cdot)$; by \eqref{eq:monic-root-bound} it is small when $z'$ is small. Shrinking $\Delta'$, we may therefore assume that these points lie in $\Delta'\times\Delta''$. They belong to $V(J)$, since $\delta^mf\in G$ for every generator $f$ of $J$ by \eqref{eq:analytic-discriminant-division}, and $\delta(z')\neq0$. As $G\subset J$, we conclude that over $\Delta'\setminus\{\delta=0\}$ the set $V(J)\cap(\Delta'\times\Delta'')$ is the union of the $q$ disjoint local holomorphic graphs just described. Thus $A_\delta$ is a complex submanifold of dimension $d$, open in $V(J)\cap(\Delta'\times\Delta'')$, and a $q$-sheeted covering of $\Delta'\setminus\{\delta=0\}$, which is connected by \cref{lem:bounded-holomorphic-removal}.

Now let $f\in\mathcal O_n$ vanish on $V(J)$. By \eqref{eq:analytic-discriminant-division}, the polynomial $R_f(z',\cdot)$ vanishes at the $q$ distinct roots $\tau_i$ for every $z'$ with $\delta(z')\neq0$ near $0$, because $\delta^mf$ and the elements of $G$ vanish at the corresponding points of $V(J)$. Its degree is less than $q$, so $R_f=0$. Hence $\delta^mf\in G\subset J$, and since $\delta\notin J$ and $J$ is prime, $f\in J$. This proves $I_{V(J),0}=J$.

\Needspace{7\baselineskip}
\emph{Step 4: density of the covering locus and the bound on the fibers.} For a complex linear form $\lambda$ on $\mathbb C^{n-d}$ and $z'\in\Delta'\setminus\{\delta=0\}$, put
\[
 P_\lambda(z',T)=\prod_{z\in\pi^{-1}(z')\cap A_\delta}
                     (T-\lambda(z'')).
\]
Its coefficients are elementary symmetric functions of the values of $\lambda$ on the sheets, hence holomorphic and bounded on $\Delta'\setminus\{\delta=0\}$; by \cref{lem:bounded-holomorphic-removal} they extend holomorphically to $\Delta'$. The function $\delta(z')P_\lambda(z',\lambda(z''))$ vanishes on $V(J)\cap(\Delta'\times\Delta'')$: on $A_\delta$ the second factor vanishes, and elsewhere the first. By Step 3 its germ belongs to $J$, and since $\delta\notin J$, the germ of $P_\lambda(z',\lambda(z''))$ belongs to $J$; so this function vanishes on $V(J)$ near $0$.

Let $z'_0\in\Delta'$ and let $F$ be the fiber over $z'_0$ of the closure $\overline{A_\delta}$ of $A_\delta$ in $\Delta'\times\Delta''$. We claim that the roots of $P_\lambda(z'_0,\cdot)$ are exactly the values $\lambda(y'')$, $y\in F$. Indeed, every $y\in F$ is a limit of points $y_j\in A_\delta$ with $\pi(y_j)\to z'_0$, and $\lambda(y_j'')$ is a root of $P_\lambda(\pi(y_j),\cdot)$; by continuity of the coefficients, $\lambda(y'')$ is a root of $P_\lambda(z'_0,\cdot)$. Conversely, choose $z'_j\to z'_0$ with $\delta(z'_j)\neq0$. The $q$ points of $A_\delta$ over $z'_j$ lie in a compact subset of $\Delta'\times\Delta''$ by Step 1, so after passing to a subsequence they converge to points of $F$; their values of $\lambda$, which are the $q$ roots of $P_\lambda(z'_j,\cdot)$, converge to the $q$ roots of $P_\lambda(z'_0,\cdot)$, counted with multiplicity. It follows that $F$ has at most $q$ points: a linear form taking distinct values at $q+1$ points of $F$ would give $q+1$ distinct roots of a polynomial of degree $q$.

It remains to show that $V(J)$ is contained in $\overline{A_\delta}$ near $0$; after shrinking the polydiscs once more, this gives the density of $A_\delta$ and the bound $q$ for all fibers. Suppose not. Then there are points $y_j\in V(J)\setminus\overline{A_\delta}$ with $y_j\to0$. For each $j$, the forms $\lambda$ with $\lambda(y_j'')=\lambda(w'')$ for some point $w$ of the finite fiber $F_j$ of $\overline{A_\delta}$ over $\pi(y_j)$ form finitely many proper linear subspaces. We choose one form $\lambda$ outside all these countably many subspaces: for instance, $\lambda$ with coefficients $(1,t,\ldots,t^{n-d-1})$, where $t$ avoids the countably many roots of the nonzero polynomials $\sum_kv_kt^{k-1}$ attached to the vectors $v=y_j''-w''$. Then $\lambda(y_j'')$ is not a root of $P_\lambda(\pi(y_j),\cdot)$ for any $j$, by the previous paragraph. On the other hand, $P_\lambda(z',\lambda(z''))$ vanishes on $V(J)$ near $0$, hence at $y_j$ for all large $j$, a contradiction.

\emph{Step 5: connectedness of the covering locus.} The covering $A_\delta$ of the connected manifold $\Delta'\setminus\{\delta=0\}$ has finitely many connected components, each of which is itself a covering with a positive number of sheets. Suppose that there are at least two of them. For each component $A^{(i)}$ and each linear form $\lambda$, the product $P^{(i)}_\lambda$ of the factors $T-\lambda(z'')$ over the points of $A^{(i)}$ again has bounded holomorphic coefficients on $\Delta'\setminus\{\delta=0\}$, which extend to $\Delta'$, and $P_\lambda=\prod_iP^{(i)}_\lambda$. Choose points $z'_j\to0$ with $\delta(z'_j)\neq0$, and, by the countable choice of Step 4, a form $\lambda$ that takes distinct values at the $q$ points of $A_\delta$ over each $z'_j$. Since the germ of $P_\lambda(z',\lambda(z''))$ belongs to the prime ideal $J$, so does the germ of some factor $P^{(i)}_\lambda(z',\lambda(z''))$, which then vanishes on $V(J)$ near $0$. Choose $k\neq i$. For large $j$ the points of $A^{(k)}$ over $z'_j$ are close to $0$, by Step 1, so $P^{(i)}_\lambda(z'_j,\cdot)$ vanishes at their values of $\lambda$; but these values differ from the roots of $P^{(i)}_\lambda(z'_j,\cdot)$, which are the values of $\lambda$ at the points of $A^{(i)}$ over $z'_j$. This contradiction shows that $A_\delta$ is connected, and completes the proof.
\end{proof}
\index{analytic subset!local parametrization}\index{branched covering!analytic parametrization}

We record a consequence that will be used repeatedly. Let $(A,0)$ be irreducible and let $(B,0)$ be an analytic subgerm different from $(A,0)$. Then for representatives on every sufficiently small parametrization polydisc $\Delta'\times\Delta''$ as in \cref{thm:analytic-local-parametrization}, the set $A\setminus B$ is dense in $A$; in particular $A_{\mathrm{reg}}\supset A_\delta$ is dense in $A$. Indeed, $B\cap A_\delta$ is an analytic subset of the connected manifold $A_\delta$. It is a proper subset, for otherwise the closed set $B$ would contain the closure $A\cap(\Delta'\times\Delta'')$ of $A_\delta$, and the two germs would agree. By \cref{lem:bounded-holomorphic-removal}, $A_\delta\setminus B$ is dense in $A_\delta$, and hence in $A$.

\begin{corollary}[Analytic Nullstellensatz]\label{cor:analytic-nullstellensatz}
For every ideal $I\subset\mathcal O_n$, the ideal of germs vanishing on $V(I)$ is the radical $\sqrt I=\{f:f^k\in I\text{ for some }k\ge1\}$.
\end{corollary}
\begin{proof}
Every element of $\sqrt I$ vanishes on $V(I)$. For the converse we first show that $\sqrt I$ is the intersection of the prime ideals containing $I$. A prime ideal containing $I$ contains $\sqrt I$. Let $f\notin\sqrt I$, so that the ideals containing $I$ and disjoint from $S=\{1,f,f^2,\ldots\}$ form a nonempty family; by \cref{prop:analytic-noetherian} it has a maximal element $\mathfrak p$, since otherwise there would be a strictly increasing infinite sequence of ideals. The ideal $\mathfrak p$ is prime: if $a,b\notin\mathfrak p$, then $\mathfrak p+(a)$ and $\mathfrak p+(b)$ meet $S$, say $f^k=p_1+xa$ and $f^l=p_2+yb$ with $p_1,p_2\in\mathfrak p$; multiplying, we see that $ab\in\mathfrak p$ would give $f^{k+l}\in\mathfrak p$, which is impossible. Thus $f$ lies outside some prime ideal containing $I$. Now let $g$ vanish on $V(I)$. For every prime ideal $J\supset I$ we have $V(J)\subset V(I)$, so $g$ vanishes on $V(J)$ and therefore $g\in J$, by the equality $I_{V(J),0}=J$ proved in \cref{thm:analytic-local-parametrization}. Hence $g$ lies in the intersection of all these prime ideals, which is $\sqrt I$.
\end{proof}
\index{Nullstellensatz!analytic}

\Needspace{10\baselineskip}
\begin{proposition}\label{prop:analytic-dimension-drop}
Let $(A,0)$ be an irreducible analytic germ. The regular points of every sufficiently small representative of $(A,0)$ all have the same dimension $d$, equal to the base dimension in \cref{thm:analytic-local-parametrization}; we call $d$ the \emph{dimension} of $(A,0)$. Every sufficiently small representative is pure-dimensional: all its local irreducible branches, at all of its points, have dimension $d$. A proper analytic subgerm of an irreducible germ has strictly smaller dimension. For a reducible germ we define the dimension as the maximum of the dimensions of its branches.
\end{proposition}
\begin{proof}
Let $A$ be the representative on a parametrization polydisc. Two local parametrizations of $(A,0)$ have covering loci which are open and dense in $A$ near $0$; their intersection is therefore nonempty, and at its points $A$ is a manifold whose dimension equals both base dimensions. Hence $d$ is determined by the germ. A regular point of $A$ has a neighborhood in $A$ which is a manifold; this neighborhood meets the dense covering locus, which has dimension $d$, so the regular point has dimension $d$ as well. Now let $y\in A$, let $B_1,\ldots,B_p$ be representatives of the branches of $(A,y)$, and fix $i$. Since $B_i$ is not contained in $\bigcup_{k\neq i}B_k$, the density property following \cref{thm:analytic-local-parametrization} gives points of the covering locus of a parametrization of $B_i$ outside the other branches. Near such a point, $A$ coincides with $B_i$ and is a manifold of dimension $\dim B_i$; such points form a nonempty open subset of $A$, which meets the covering locus of $A$. Hence $\dim B_i=d$.

Let $(B,0)\subset(A,0)$ be germs with both irreducible, and let $J_A\subset J_B$ be their prime ideals. In Step 1 of the proof of \cref{thm:analytic-local-parametrization}, choose each direction so that it serves for a nonzero element of $J_B\cap\mathcal O_k$ and also, as long as $J_A\cap\mathcal O_k\neq0$, for a nonzero element of $J_A\cap\mathcal O_k$; this gives parametrizations of both germs in common coordinates. Since $J_A\subset J_B$, the ideal $J_B\cap\mathcal O_k$ is nonzero whenever $J_A\cap\mathcal O_k$ is, so $\dim B\leq\dim A$. Suppose that the dimensions are equal, so that both germs are finite over the same base, and let $\delta_A$, $\delta_B$ be the two discriminants. By \cref{lem:bounded-holomorphic-removal}, the set $A'$ of points of the connected covering locus of $A$ lying over $\{\delta_B\neq0\}$ is a connected manifold of dimension $d$, dense in $A$. The points of $B$ over $\{\delta_A\delta_B\neq0\}$ form a nonempty submanifold of $A'$ of the same dimension, hence an open subset of $A'$. Thus $B\cap A'$ is an analytic subset of the connected manifold $A'$ with nonempty interior, and $B\supset A'$ by \cref{lem:bounded-holomorphic-removal}. Since $B$ is closed and $A'$ is dense in $A$, we get $(B,0)=(A,0)$. Hence a proper subgerm has strictly smaller dimension. For a general proper subgerm $B$ of $(A,0)$, each branch of $B$ is a proper subgerm of $(A,0)$, and the same conclusion follows.
\end{proof}
\index{analytic subset!dimension}

The dimension of an analytic set $E$ at a point $x$ is the dimension of its germ $(E,x)$. Dimension is monotone: if $(E,x)\subset(F,x)$, every branch of $(E,x)$ is contained in a branch of $(F,x)$, so $\dim(E,x)\leq\dim(F,x)$. It is also upper semicontinuous: since small representatives of the branches of $(E,x)$ are pure-dimensional, every branch of $(E,y)$, for $y$ near $x$, is contained in one of them and has dimension at most $\dim(E,x)$. We say that $E$ is \emph{pure-dimensional} if its dimension is the same at all points, and that $E\subset M$ has \emph{codimension at least $k$} if $\dim(E,x)\leq\dim_xM-k$ for all $x\in E$.

\Needspace{9\baselineskip}
\begin{theorem}[Cartan's coherence theorem]\label{thm:analytic-ideal-coherence}
For an analytic subset $A$ of a complex manifold $M$, the sheaf $\mathcal I_A$ of holomorphic functions vanishing on $A$ is coherent.
\end{theorem}
\begin{proof}
The assertion is local. Near a point of $A$, the set $A$ is the union of representatives $A_1,\ldots,A_p$ of the branches of its germ, and $\mathcal I_A=\bigcap_i\mathcal I_{A_i}$. The intersection of two coherent subsheaves $\mathcal F,\mathcal G$ of $\mathcal O_M$ is the kernel of the morphism $\mathcal F\oplus\mathcal G\to\mathcal O_M$, $(a,b)\mapsto a-b$, and is coherent by \cref{thm:oka-coherence,lem:coherent-sheaf-algebra}. It therefore suffices to prove that $\mathcal I_A$ is coherent near $0$ when $(A,0)\subset(\mathbb C^n,0)$ is irreducible. We use the parametrization of \cref{thm:analytic-local-parametrization}, with its primitive form $u$ and discriminant $\delta$, and the exponent $m$ of \eqref{eq:analytic-discriminant-division}. Note that $A$ may have several branches at points near $0$; this is why we need an argument that works uniformly on a neighborhood.

\emph{Step 1: dependence on the primitive form.} For $c=(c_{d+1},\ldots,c_n)\in\mathbb C^{n-d}$ put $u_c=\sum_kc_kz_k$. For $z'\in\Delta'\setminus\{\delta=0\}$ let $z^{(1)},\ldots,z^{(q)}$ be the points of $A$ over $z'$, and define
\[
 W_c(z',T)=\prod_{i=1}^q\bigl(T-u_c(z^{(i)})\bigr),\qquad
 \delta_c(z')=\prod_{i<l}\bigl(u_c(z^{(i)})-u_c(z^{(l)})\bigr)^2,\qquad
 B_{kc}(z',T)=\sum_{j<q}D_c(z')D_{kcj}(z')T^j,
\]
where $D_c$ is the Vandermonde determinant $\det(u_c(z^{(i)})^j)_{i,j}$ and $D_{kcj}$ is obtained from it by replacing the $j$th column with the coordinates $z_k^{(i)}$. These expressions do not change under permutations of the points $z^{(i)}$, since $D_c$ and $D_{kcj}$ both change sign under a transposition. Hence their coefficients in $T$ are polynomials in $c$ whose coefficients are bounded holomorphic functions on $\Delta'\setminus\{\delta=0\}$, and these extend holomorphically to $\Delta'$ by \cref{lem:bounded-holomorphic-removal}. On $A_\delta$ we have $W_c(z',u_c(z''))=0$ and, by Cramer's rule in the form $\sum_jD_{kcj}u_c(z^{(i)})^j=D_cz_k^{(i)}$, also $\delta_c(z')z_k=B_{kc}(z',u_c(z''))$, for every $c$. Let $G_1,\ldots,G_N$ be the coefficients, as polynomials in $c$, of the functions $W_c(z',u_c)$ and $\delta_cz_k-B_{kc}(z',u_c)$, $d<k\leq n$. They are holomorphic on $\Delta'\times\Delta''$ and vanish on $A_\delta$, hence on $A$, which is the closure of $A_\delta$. For the primitive form $u$ of the parametrization, $\delta_c\neq0$ off $\{\delta=0\}$; so the polynomial $\delta_c^m=\sum_\alpha h_\alpha(z')c^\alpha$ in $c$ is not identically zero, and we choose a coefficient $h=h_\alpha\not\equiv0$.

\emph{Step 2: a description of the stalks.} We claim that on a small polydisc around $0$,
\begin{equation}\label{eq:analytic-coherence-colon}
 \mathcal I_{A,x}=\{f\in\mathcal O_{M,x}:h f\in(G_{1,x},\ldots,G_{N,x})\}
\end{equation}
for every point $x$. If $hf\in(G_{1,x},\ldots,G_{N,x})$, then $hf$ vanishes on $A$ near $x$, so $f$ vanishes on $A\setminus\{h=0\}$ near $x$. This set is dense in $A$: the covering locus $A_\delta$ is dense in $A$, and since $\pi$ is a local biholomorphism on $A_\delta$, the set where $h\circ\pi$ vanishes is nowhere dense in $A_\delta$. Hence $f\in\mathcal I_{A,x}$, which proves the inclusion from right to left.

\emph{Step 3: the inclusion from left to right at points of $A$.} Fix $x\in A$ and $f\in\mathcal I_{A,x}$. Call $c$ admissible if $\delta_c\not\equiv0$ and $u_c$ takes distinct values at the points of the finite fiber of $A$ over $x'$; the admissible $c$ form a dense open subset of $\mathbb C^{n-d}$, the complement of an algebraic subset and finitely many proper linear subspaces. Fix an admissible $c$. By Step 4 of the proof of \cref{thm:analytic-local-parametrization}, the roots of $W_c(x',\cdot)$ are the values of $u_c$ at the points of the fiber over $x'$. By the preparation theorem, as in the proof of \cref{thm:oka-coherence}, $W_c=W_c^xQ_c^x$ near $(x',u_c(x''))$, where $W^x_c$ is a Weierstrass polynomial in $T-u_c(x'')$ and $Q^x_c$ is a polynomial in $T$ with $Q^x_c(x',u_c(x''))\neq0$. For $z'$ near $x'$ with $\delta(z')\delta_c(z')\neq0$, the roots of $W^x_c(z',\cdot)$ are distinct and are exactly the values of $u_c$ at the points of $A$ over $z'$ that are close to $x$; let $q_x$ be their number. Similarly, the coordinate $z_k$ of each point of $A$ is a root of the minimal polynomial $W_k(z',\cdot)$ of Step 2 of that proof, and we let $W^x_k$ be the Weierstrass factor of $W_k$ at $(x',x_k)$; it vanishes at the coordinates $z_k$ of the points of $A$ close to $x$, and its degree is at most $q$.

Divide $f$ successively by the Weierstrass polynomials $W^x_n(z',z_n),\ldots,W^x_{d+1}(z',z_{d+1})$ at $x$. As in the proof of \eqref{eq:analytic-discriminant-division}, substituting $\delta_cz_k\equiv B_{kc}(z',u_c)$ and dividing by $W^x_c$, we obtain
\[
 \delta_c^mf\equiv R(z',u_c)\pmod{(W_c(z',u_c),\ \delta_cz_k-B_{kc}(z',u_c))_x},
\]
with a polynomial $R$ of degree less than $q_x$. Here we used that $W^x_c(z',u_c)$ is a unit at $x$ times $W_c(z',u_c)$, and that $\delta_c^qW^x_k(z',z_k)$ belongs to the ideal on the right. For the latter, write $\delta_c^{q_k}W^x_k(z',z_k)\equiv H(z',u_c)$, with $q_k=\deg W^x_k\leq q$ and $H(z',T)$ obtained by substituting $B_{kc}(z',T)/\delta_c$ for $z_k$ and clearing denominators. At the $q_x$ distinct roots of $W^x_c(z',\cdot)$, for generic $z'$ near $x'$, the polynomial $H(z',\cdot)$ takes the values $\delta_c^{q_k}W^x_k(z',z_k)$ at points of $A$ close to $x$, which vanish; so the remainder of $H$ after division by $W^x_c$ is zero, and $H(z',u_c)$ is a multiple of $W^x_c(z',u_c)$. Now $f$ and the elements of the ideal on the right vanish on $A$ near $x$; so $R(z',\cdot)$ vanishes at the $q_x$ distinct roots of $W^x_c(z',\cdot)$ for generic $z'$ near $x'$, and $R=0$. Since $W_c(z',u_c)$ and $\delta_cz_k-B_{kc}(z',u_c)$ are complex linear combinations of $G_1,\ldots,G_N$, we have shown that
\[
 \delta_c^m f\in(G_{1,x},\ldots,G_{N,x})\qquad\text{for every admissible }c.
\]
The polynomials in $c$ of degree at most $\deg_c\delta_c^m$ form a finite-dimensional space, and the evaluations at points of the dense open set of admissible $c$ span its dual, since a polynomial vanishing on a nonempty open set is zero. Hence each coefficient $h_\alpha$ of $\delta_c^m$ is a complex linear combination of finitely many values $\delta_{c^{(l)}}^m$ with admissible $c^{(l)}$, and $h_\alpha f\in(G_{1,x},\ldots,G_{N,x})$. This proves \eqref{eq:analytic-coherence-colon} at points of $A$.

\emph{Step 4: points outside $A$, and conclusion.} Let $x\notin A$. Choose $c$ such that $u_c(x'')$ differs from the values of $u_c$ at the finitely many points of $A$ over $x'$. Then $W_c(x',u_c(x''))\neq0$, so one of the coefficients $G_l$ does not vanish at $x$, and both sides of \eqref{eq:analytic-coherence-colon} are the unit ideal. The right side of \eqref{eq:analytic-coherence-colon} is the projection to the first coordinate of the kernel of the row matrix $(h,-G_1,\ldots,-G_N)$, which is coherent by \cref{thm:oka-coherence,lem:coherent-sheaf-algebra}. Hence $\mathcal I_A$ is coherent.
\end{proof}
\index{Cartan coherence theorem}

\begin{corollary}\label{cor:analytic-strata-hausdorff}
Let $A$ be an analytic subset of a complex manifold. Its singular locus $A_{\mathrm{sing}}$ is an analytic subset, and at every point of $A$ the intersection of $A_{\mathrm{sing}}$ with each irreducible branch of $A$ has smaller dimension than that branch. An analytic set of dimension at most $d$ at every point is locally a countable union of real submanifolds of dimensions at most $2d$. In particular, its Hausdorff $s$-measure is zero for every $s>2d$.
\end{corollary}
\begin{proof}
We first describe the ideal of a submanifold. If $N=\{w_1=\cdots=w_k=0\}$ in local coordinates $(v,w)$, then every holomorphic function $g$ vanishing on $N$ near a point satisfies $g(v,w)=\sum_lw_l\int_0^1\partial_{w_l}g(v,tw)\,\mathrm dt$, so its ideal is generated by $w_1,\ldots,w_k$. The quotient of $\mathcal O_n$ by this ideal is a ring of germs in the variables $v$, a domain; so the germ of a connected submanifold is irreducible by \cref{prop:analytic-irreducible-germs}, and its dimension is the dimension of the manifold.

Let $A$ be a small representative of an irreducible germ of dimension $d$ in $\mathbb C^n$, pure-dimensional by \cref{prop:analytic-dimension-drop}, and choose generators $g_1,\ldots,g_N$ of $\mathcal I_A$ on a neighborhood, by \cref{thm:analytic-ideal-coherence}. At a regular point, these generate the ideal of the manifold $A$, whose generators are $n-d$ normal coordinates; differentiating, we see that the matrix of differentials of the $g_j$ has rank $n-d$ there. Every minor of size $n-d+1$ therefore vanishes on the dense covering locus, hence on $A$, so the rank is at most $n-d$ everywhere on $A$. Conversely, if it equals $n-d$ at $y\in A$, choose $n-d$ of the $g_j$ with independent differentials at $y$; by the implicit function theorem their common zero set is a $d$-dimensional submanifold $N$ near $y$ containing $A$. Every branch of $(A,y)$ has dimension $d$ and is contained in the irreducible germ $(N,y)$ of dimension $d$, so it equals $(N,y)$ by \cref{prop:analytic-dimension-drop}; hence $A$ is a manifold near $y$. Thus $A_{\mathrm{sing}}$ is the subset of $A$ where all minors of size $n-d$ of this matrix vanish, an analytic subset. At every point $y$ of $A$, its intersection with each branch $B$ of $(A,y)$ is a subgerm of $B$ different from $B$, because $B$ contains regular points of $A$ arbitrarily close to $y$, namely points of the covering locus of a parametrization of $B$ outside the other branches of $(A,y)$; hence it has smaller dimension than $B$, by \cref{prop:analytic-dimension-drop}.

For a general analytic set $A$ near a point $x$, let $A_1,\ldots,A_p$ be small representatives of the branches of $(A,x)$. A point $y$ near $x$ lying on only one $A_i$ is regular for $A$ if and only if it is regular for $A_i$. A point $y\in A_i\cap A_j$, $i\neq j$, is singular. Indeed, otherwise $(A,y)$ would be irreducible, hence equal to some $(A_k,y)$, and we may assume $k\neq i$; then $(A_i,y)\subset(A_i\cap A_k,y)$. But $\dim(A_i\cap A_k,x)<\dim A_i$ by \cref{prop:analytic-dimension-drop}, so $\dim(A_i\cap A_k,y)<\dim A_i=\dim(A_i,y)$ by upper semicontinuity and purity, contradicting monotonicity. Hence near $x$
\[
 A_{\mathrm{sing}}=\bigcup_i(A_i)_{\mathrm{sing}}\cup\bigcup_{i<j}(A_i\cap A_j),
\]
an analytic subset, and its intersection with each branch has smaller dimension than the branch, by the irreducible case and by \cref{prop:analytic-dimension-drop} applied to the intersections.

Now let $A$ have dimension at most $d$ at every point. Then $A=A_{\mathrm{reg}}\cup A_{\mathrm{sing}}$, where $A_{\mathrm{reg}}$ is a complex submanifold of dimension at most $d$, and $A_{\mathrm{sing}}$ is analytic of smaller dimension near each point. Repeating this decomposition with $A_{\mathrm{sing}}$, and so on, we reach the empty set near each point after at most $d+1$ steps. Each submanifold obtained is a countable union of coordinate charts, so $A$ is locally a countable union of real submanifolds of dimensions at most $2d$. Each such chart is a countable union of images of compact cubes in $\mathbb R^r$, $r\leq2d$, under Lipschitz maps. Such an image has zero Hausdorff $s$-measure for $s>r$: covering the cube by $O(\delta^{-r})$ cubes of side $\delta$, whose images have diameter $O(\delta)$, we bound its $s$-dimensional Hausdorff premeasure at scale $O(\delta)$ by $O(\delta^{s-r})$, which tends to zero. Countable subadditivity completes the proof.
\end{proof}
\index{analytic subset!singular locus}\index{Hausdorff measure!analytic sets}

An analytic subset $Z$ of a complex manifold is \emph{irreducible} if it is nonempty and is not the union of two analytic subsets different from $Z$.

\begin{corollary}\label{cor:analytic-global-components}
Let $A$ be an analytic subset of a complex manifold $M$. The regular locus $A_{\mathrm{reg}}$ is dense in $A$. The closures of the connected components of $A_{\mathrm{reg}}$ are irreducible analytic subsets of $A$, called the \emph{irreducible components} of $A$; they form a locally finite family with union $A$, and none of them is contained in the union of the others. Every locally finite family of irreducible analytic subsets with union $A$, none of which is contained in another, consists of these components. If $M$ has a countable basis, there are at most countably many irreducible components. The regular locus of an irreducible component $Z$ is connected, and a holomorphic function defined near $Z$ that vanishes on a nonempty relatively open subset of $Z_{\mathrm{reg}}$ vanishes on $Z$.
\end{corollary}
\begin{proof}
Fix $x\in A$, let $A_1,\ldots,A_p$ be representatives of the branches of $(A,x)$, and for each $i$ choose a parametrization polydisc $P_i$ of $(A_i,x)$ so small that $A\cap P_i=\bigcup_j A_j\cap P_i$. By \cref{cor:analytic-strata-hausdorff}, $A_{\mathrm{sing}}\cap A_i$ is an analytic subset of $A_i$ whose germ at $x$ differs from $(A_i,x)$. Hence $A_{\mathrm{sing}}$ meets the connected covering locus of $A_i$ in a proper analytic subset, and by \cref{lem:bounded-holomorphic-removal} and the density property following \cref{thm:analytic-local-parametrization}, the covering locus minus $A_{\mathrm{sing}}$ is connected and dense in $A_i\cap P_i$. The set $R_i=A_i\cap P_i\cap A_{\mathrm{reg}}$ lies between this connected set and its closure, so it is connected and dense in $A_i\cap P_i$. In particular $A_{\mathrm{reg}}$ is dense in $A$.

Let $\{S_\alpha\}$ be the connected components of the manifold $A_{\mathrm{reg}}$, and $Z_\alpha=\overline{S_\alpha}$. Each $R_i$ above is connected, so it lies in a single component $S_{\alpha(i)}$. Let $U=\bigcap_iP_i$. Then $S_\alpha\cap U$ is contained in the union of the $A_i$ with $\alpha(i)=\alpha$, because points of $A_{\mathrm{reg}}\cap A_i\cap U$ lie in $R_i$; conversely $A_i\cap U\subset\overline{R_i}\subset Z_{\alpha(i)}$. Hence
\[
 Z_\alpha\cap U=\bigcup_{i:\,\alpha(i)=\alpha}A_i\cap U.
\]
Thus each $Z_\alpha$ is analytic, only finitely many of them meet $U$, and their union is $A$. Since $S_\alpha$ is open in $A$, a point of $S_\alpha$ has a neighborhood meeting no other $S_\beta$, hence not in $Z_\beta$ for $\beta\neq\alpha$; so $S_\alpha\cap Z_\beta=\varnothing$ and no $Z_\alpha$ is contained in the union of the others. The same remark shows that $S_\alpha\subset(Z_\alpha)_{\mathrm{reg}}$: near a point of $S_\alpha$, the set $Z_\alpha$ coincides with the manifold $A$. Hence $(Z_\alpha)_{\mathrm{reg}}$ lies between the connected set $S_\alpha$ and its closure, and is connected. If $Z_\alpha$ were the union of two analytic subsets $Z'$, $Z''$ different from it, then $S_\alpha\cap Z'$ and $S_\alpha\cap Z''$ would be analytic subsets of the connected manifold $S_\alpha$ different from $S_\alpha$, since otherwise their closure would be $Z_\alpha$; by \cref{lem:bounded-holomorphic-removal} they would be closed and nowhere dense in $S_\alpha$, and so would their union, which is $S_\alpha$. This contradiction shows that $Z_\alpha$ is irreducible.

If $M$ has a countable basis, so does $A_{\mathrm{reg}}$; each $S_\alpha$ is open in $A_{\mathrm{reg}}$ and contains a basic open set, and disjoint components contain different ones, so there are at most countably many components. Finally, a holomorphic function vanishing on a nonempty open subset of the connected manifold $(Z_\alpha)_{\mathrm{reg}}$ vanishes on all of it by the identity theorem, and hence on its closure $Z_\alpha$.

It remains to prove uniqueness. Let $A=\bigcup_\beta Y_\beta$ be a locally finite decomposition into irreducible analytic subsets, none contained in another. A connected manifold is not a locally finite union of closed nowhere dense subsets, since such a union is closed and nowhere dense. Hence, by \cref{lem:bounded-holomorphic-removal}, each $S_\alpha$ lies in some $Y_\beta$, and then $Z_\alpha\subset Y_\beta$. Conversely, fix $\beta$. By what we have proved, applied to the analytic set $Y_\beta$, the set $Y_\beta$ is the union of its irreducible components, and it has only one, for otherwise it would be the union of one of them and of the union of the others, two analytic subsets different from $Y_\beta$. Hence $(Y_\beta)_{\mathrm{reg}}$ is a connected manifold, dense in $Y_\beta$ and covered by the locally finite family of closed analytic subsets $(Y_\beta)_{\mathrm{reg}}\cap Z_\alpha$; as before, $Y_\beta\subset Z_\alpha$ for some $\alpha$. Choosing $\beta'$ with $Z_\alpha\subset Y_{\beta'}$, we get $Y_\beta\subset Y_{\beta'}$, so $\beta'=\beta$ and $Y_\beta=Z_\alpha$. Thus every $Y_\beta$ is an irreducible component of $A$; and every component $Z_\alpha$, being contained in some $Y_\beta=Z_{\alpha'}$, equals it, because $S_\alpha\cap Z_{\alpha'}=\varnothing$ for $\alpha'\neq\alpha$.
\end{proof}
\index{analytic subset!global components}

\begin{corollary}\label{cor:analytic-codimension-two-psh-removal}
Let $A$ be an analytic subset of a complex manifold $M$ of codimension at least two. Then every plurisubharmonic function $u$ on $M\setminus A$ is locally bounded above near $A$ and extends uniquely to a plurisubharmonic function on $M$, whose value at $a\in A$ is $\limsup_{M\setminus A\ni y\to a}u(y)$.
\end{corollary}
\begin{proof}
The assertion is local, and extensions defined on overlapping charts agree, since their values on $A$ are given by the same formula. Near a point where $M$ has dimension at most one, the codimension condition forces $A=\varnothing$. Otherwise, let $\Omega$ be a connected coordinate chart of dimension $n\geq2$. Then $\dim A\leq n-2$, and \cref{cor:analytic-strata-hausdorff} shows that $A\cap\Omega$ is locally $\mathcal H^{2n-3}$-null, since $2n-3>2(n-2)$. By \cref{thm:psh-small-set-extension}, $u$ extends uniquely to a plurisubharmonic function on $\Omega$ with the stated values on $A$; being upper semicontinuous, the extension is locally bounded above.
\end{proof}
\index{plurisubharmonic function!codimension-two removal}

For an irreducible analytic subset $D$ of codimension one of a complex manifold $M$, we write $[D]$ for the current of integration over the complex hypersurface $D_{\mathrm{reg}}$ of $M\setminus D_{\mathrm{sing}}$ (\cref{ex:submanifold-current}), with coefficient measures extended by zero across $D_{\mathrm{sing}}$. The following proof shows that $[D]$ is a closed positive current with locally finite mass.

\begin{theorem}\label{thm:analytic-order-zero-support}
Let $A$ be a nowhere dense analytic subset of a complex manifold $M$, and let $R$ be a closed real $(1,1)$-current on $M$ with measure coefficients and support in $A$. Then
\[
 R=\sum_D c_D[D],
\]
where $D$ ranges over the irreducible components of $A$ of codimension one, the real coefficients $c_D$ are uniquely determined, and the sum is locally finite. If $R$ is positive, then every $c_D\geq0$. In particular $R=0$ if $A$ has codimension at least two.
\end{theorem}
\begin{proof}
\emph{Step 1: supports of high codimension.} Suppose first that, near a point, the support of $R$ lies in a real submanifold $\{g_1=\cdots=g_k=0\}$ with independent differentials $\mathrm dg_1,\ldots,\mathrm dg_k$, $k\geq3$. Since $R$ has measure coefficients and $g_j$ vanishes on its support, $g_jR=0$; since $R$ is closed, $\mathrm dg_j\wedge R=\mathrm d(g_jR)=0$. In real coordinates whose first three are $g_1,g_2,g_3$, a coefficient of $R$ in front of $\mathrm dx_a\wedge\mathrm dx_b$ is annihilated by wedging with $\mathrm dg_l$ only if $l\in\{a,b\}$; so every coefficient would have to involve all three indices $1,2,3$, and $R=0$ near the point. Now let $A$ have codimension at least two. Near a point of $A_{\mathrm{reg}}$, the set $A$ is a complex submanifold of real codimension at least four, so $R$ vanishes there, and the support of $R$ lies in $A_{\mathrm{sing}}$. This is analytic of smaller dimension near each point by \cref{cor:analytic-strata-hausdorff}, and again of codimension at least two; repeating the argument, we find after finitely many steps near each point that $R=0$. This proves the last assertion.

\emph{Step 2: regular points of codimension one.} Let $y\in A_{\mathrm{reg}}$ be a point where $A$ is a complex hypersurface. Choose holomorphic coordinates $(\zeta,w)$ near $y$ with $A=\{\zeta=0\}$, write $\zeta=\xi+\mathrm i\eta$, and let $t\in\mathbb R^{2n-2}$ be the real coordinates of $w$. The identities $\mathrm d\xi\wedge R=\mathrm d\eta\wedge R=0$ of Step 1 force $R=\mu\,\mathrm d\xi\wedge\mathrm d\eta$ for a real measure $\mu$ supported on $\{\xi=\eta=0\}$, so $\mu=\delta_0(\xi,\eta)\otimes\sigma(t)$ for a real measure $\sigma$. Closedness of $R$ says that all distributional derivatives $\partial\sigma/\partial t_l$ vanish. On a smaller ball in the $t$-variables, the convolutions of $\sigma$ with a smoothing kernel therefore have zero gradient and are constant; integrating them against a fixed test function of integral one, we see that these constants converge as the kernel shrinks, and that $\sigma$ is a constant multiple $c$ of Lebesgue measure. Hence $R=c[\{\zeta=0\}]$ near $y$, with $c$ real. (If $n=1$, $\sigma$ is a measure on a point and the same conclusion holds.)

\emph{Step 3: the currents $[D]$.} Let $D$ be an irreducible analytic subset of $M$ of codimension one. By the uniqueness in \cref{cor:analytic-global-components}, applied to $A=D$, the set $D$ is its own only irreducible component; so $D_{\mathrm{reg}}$ is a connected manifold, dense in $D$, and its dimension is $n-1$. Every branch of $D$ at every point contains regular points of $D$, so all these branches have dimension $n-1$. Let $x\in D$ and let $D_1,\ldots,D_k$ be representatives of the branches of $(D,x)$. In the parametrization of $(D_i,x)$ we have $d=n-1$, and the minimal polynomial $W_n(z',T)$ of the class of $z_n$ in Step 2 of the proof of \cref{thm:analytic-local-parametrization} is a Weierstrass polynomial. It generates the prime ideal $I_{D_i,x}$: the remainder of the division of an element of this ideal by $W_n(z',z_n)$ is a polynomial in $z_n$ of degree less than $\deg W_n$ that vanishes at the class of $z_n$, hence it is zero by minimality. Let $f_i$ be a holomorphic function near $x$ with this germ. The sheaf $\mathcal I_{D_i}$ has finitely many generators near $x$ by \cref{thm:analytic-ideal-coherence}, and each of them is a multiple of $f_i$ at $x$, hence on a neighborhood; so $f_i$ generates $\mathcal I_{D_i,y}$ for $y$ near $x$. Let $y\in D_{\mathrm{reg}}$ be close to $x$. By the proof of \cref{cor:analytic-strata-hausdorff}, $y$ lies on only one $D_i$, and $D=D_i=\{\tau=0\}$ near $y$ for a local coordinate $\tau$; then $f_i$ is a unit times $\tau$, and $f_j(y)\neq0$ for $j\neq i$. Thus $f=f_1\cdots f_k$ vanishes to order one along $D_{\mathrm{reg}}$ near $x$, and $f^{-1}(0)=D$ near $x$. In the notation of \cref{sec:poincare-lelong}, the regular part of the zero set of $f$ is $D_{\mathrm{reg}}$ near $x$, with multiplicity one. By \cref{thm:poincare-lelong}, $[D]=\ddc\log|f|^2$ near $x$; hence $[D]$ is closed and positive, with locally finite mass.

\emph{Step 4: conclusion.} Let $D=Z_\alpha$ be an irreducible component of codimension one, the closure of the connected component $S_\alpha$ of $A_{\mathrm{reg}}$. By Step 2, near each point of $S_\alpha$ we have $R=c[D]$ for a real constant $c$, and this constant is locally constant on the connected manifold $S_\alpha$; we call it $c_D$. By \cref{cor:analytic-global-components} the sum $\sum_Dc_D[D]$ is locally finite, and by Step 3 it is a closed real current with measure coefficients. The difference $R'=R-\sum_Dc_D[D]$ vanishes near every point of $A_{\mathrm{reg}}$: near a point of $S_\alpha$ with $Z_\alpha$ of codimension one by Step 2, since the other components do not meet $S_\alpha$; and near a point of $A_{\mathrm{reg}}$ where $A$ has codimension at least two by Step 1, since no component of codimension one passes through such a point. Thus the support of $R'$ lies in $A_{\mathrm{sing}}$, which has codimension at least two because $A$ is nowhere dense and $A_{\mathrm{sing}}$ has smaller dimension than each branch of $A$. By Step 1, $R'=0$. The coefficients are unique, because near a point of $S_\alpha$ the current $R$ determines $c_D$; and if $R\geq0$, then $c_D\geq0$, because $[D]$ is positive and nonzero near such a point.
\end{proof}
\index{current!analytic support theorem}

\section{Reduced spaces and normalization}
\label{sec:reduced-spaces-normalization}

A \emph{reduced complex space} is a Hausdorff space locally identified with an analytic subset $A$ of a domain in $\mathbb C^N$, equipped with the sheaf
$\mathcal O_A=(\mathcal O/\mathcal I_A)|_A$; the local identifications preserve this sheaf. Thus its holomorphic functions locally extend to the ambient domain. We use a countable basis; compact spaces with such local models automatically have one, by a finite chart cover. A holomorphic map is locally a map into ambient coordinates whose components belong to the source structure sheaf. Its coordinate functions consequently extend near each source point, and composition is holomorphic. The regular and singular loci, branches and dimensions are intrinsic, since local holomorphic isomorphisms identify the regular manifold parts and their dense parametrization loci. All complex spaces below are reduced.
\index{complex space!reduced}

The structure sheaf is coherent: in a chart it is the quotient by the coherent ideal of \cref{thm:analytic-ideal-coherence}. The sheaf algebra of \cref{lem:coherent-sheaf-algebra} applies in these charts also to $\mathcal O_X$-modules. A local ring $\mathcal O_{X,x}$ is a quotient of $\mathcal O_N$, so it is Noetherian. Coherent ideal sheaves define analytic subspaces by their zero sets; for a reduced analytic subset use its vanishing ideal. A \emph{finite} holomorphic map means a proper map with finite fibers, and \emph{bimeromorphic} means that the map is an isomorphism off nowhere dense analytic subsets of both spaces. Properness means compact inverse images of compact sets. Such maps between these locally compact spaces are closed: for a convergent sequence in the image of a closed set, lift it into the compact inverse image of a compact neighborhood and pass to a convergent subsequence.
\index{holomorphic map!finite}\index{bimeromorphic map}

A \emph{weakly holomorphic} function on $X$ is holomorphic on $X_{\mathrm{reg}}$ and locally bounded near the singular locus. Denote its sheaf by $\widetilde{\mathcal O}_X$. On a reducible germ it is the direct sum of the corresponding sheaves on the local branches: their regular parts are disjoint outside pairwise intersections, and bounded holomorphic removal on each regular branch fills those intersections. The total ring of fractions of a reduced local ring inverts germs which are not identically zero on any branch. These are exactly the nonzero divisors. Indeed, such a germ times $g$ equal to zero forces $g=0$ by density on every branch; if it vanishes identically on a branch, multiply it by a product of equations vanishing on the other branches but not on that one to obtain a nonzero annihilated germ.
\index{holomorphic function!weakly holomorphic}

\Needspace{12\baselineskip}
\begin{proposition}\label{prop:weak-holomorphic-denominator}
Every point of $X$ has a neighborhood $U$ and a holomorphic function $h$, nonzero on every local branch, such that
\[
 h_y\widetilde{\mathcal O}_{X,y}\subset\mathcal O_{X,y}
 \quad(y\in U).
\]
Weakly holomorphic germs are precisely the integral elements in the total ring of fractions of $\mathcal O_{X,x}$, and form a finite $\mathcal O_{X,x}$-module. On an irreducible germ every weakly holomorphic function has one finite limit at its base point.
\end{proposition}
\begin{proof}
First use a local parametrization of an irreducible germ, with primitive coordinate $u$, degree $q$ and discriminant $\delta$. For a bounded weakly holomorphic germ $f$ near a point $y$, separate the finitely many fibers over $y'$. On the sheets near $y$ use the values of $f$, and on the other sheets use zero. Interpolate those $q$ values by $\sum_{j<q}b_j(z')u^j$. The coefficients are holomorphic off the discriminant, invariant under permutations within each sheet group. Cramer's rule shows that $\delta b_j$ is bounded there: its expression is a product of Vandermonde determinants with bounded sheet coordinates and bounded values of $f$. It therefore extends holomorphically on the base. The expression $\sum_{j<q}(\delta b_j)u^j$ is a holomorphic germ on $X$ and equals $\delta f$ on its dense regular covering part near $y$. Thus the same $\delta$ is a denominator at every nearby point.

For finitely many branches $X_j$, choose a denominator $\delta_j$ on each and lift it to the ambient domain. Add to it a generic linear combination of equations vanishing on $X_j$ so that it is nonzero on every other branch too; finitely many proper conditions are excluded. Then $\delta=\prod_j\delta_j$ is nonzero on all branches and works separately on each. Choose $v_j$ vanishing on every other branch and not identically on $X_j$, by products of separating equations. Set $h=\delta\sum_jv_j$. On $X_j$, $h=\delta v_j$, so it is a nonzero divisor. If $\delta f=g_j$ on the $j$th branch, with $g_j$ lifted holomorphically, then $hf=\sum_jv_jg_j$ on all branches. This gives the uniform denominator. It embeds the weak germs in $h^{-1}\mathcal O_{X,x}$; multiplication by $h$ identifies them with a submodule of the Noetherian ring $\mathcal O_{X,x}$, proving finite generation.

On one branch, the elementary symmetric functions of the $q$ values of $f$ on the regular fibers are bounded holomorphic base functions. Removal extends them, giving a monic equation for $f$ over $\mathcal O_d\subset\mathcal O_{X,x}$. For several branches take the product of their equations after lifting their coefficients. Hence every weak germ is integral. Conversely an integral meromorphic germ is bounded near $x$ by the monic root bound. On the regular locus its bounded meromorphic singularities are removable by \cref{lem:bounded-holomorphic-removal}, so it is weakly holomorphic. Finally, on an irreducible germ the regular locus has a basis of connected neighborhoods by local parametrization and density. The cluster set of a bounded holomorphic function there is a nonempty connected compact set, the intersection of its decreasing connected compact image closures. Its integral equation restricts that cluster set to finitely many roots at $x$. Thus it is a singleton, proving the finite limit.
\end{proof}
\index{holomorphic function!universal denominator}\index{integral closure!analytic local ring}

A point $x$ is \emph{normal} if its germ is irreducible and $\mathcal O_{X,x}$ is integrally closed in its fraction field. By the preceding proposition this means $\widetilde{\mathcal O}_{X,x}=\mathcal O_{X,x}$. Conversely that equality forces irreducibility: the weak function equal to one on one local branch and zero on the others could not be a continuous holomorphic function at their common point. A space is normal if every point is normal. Regular points are normal by \cref{lem:bounded-holomorphic-removal}.
\index{complex space!normal}

\Needspace{7\baselineskip}
\begin{lemma}\label{lem:normal-locus-open}
The nonnormal locus is analytic. Bounded holomorphic functions off a nowhere dense analytic subset of a normal space extend holomorphically across it.
\end{lemma}
\begin{proof}
We first record two finite-sheaf facts. The sheaf $\operatorname{Hom}(\mathcal F,\mathcal G)$ is coherent for coherent modules: using a finite presentation of $\mathcal F$, it is the kernel of a morphism between finite sums of $\mathcal G$. The support of a coherent module is analytic. In a presentation $\mathcal O^a\to\mathcal O^b\to\mathcal F\to0$, its zero stalks are exactly the points where the matrix has rank $b$ after evaluation. This uses the following elementary local-ring argument: if a finite module $M$ satisfies $M=\mathfrak m M$, write each generator as a combination with coefficients in the maximal ideal. The determinant of the resulting $I-C$ is a unit and annihilates the generators by the adjugate identity, so $M=0$. Applied to the cokernel, this proves the assertion; the rank failure is defined by minors.

Choose the uniform denominator $h$ of \cref{prop:weak-holomorphic-denominator}, and let $\mathcal I$ be the coherent vanishing ideal of $\{h=0\}$ in $X$. Put $\mathcal F=\operatorname{Hom}(\mathcal I,\mathcal I)$. Multiplication embeds $\mathcal O_X$ in $\mathcal F$. Any endomorphism is multiplication by a meromorphic germ: if its value on $h$ is $a$, linearity gives $h\Phi(g)=ag$ for $g\in\mathcal I$. The finite-module determinant argument, applied to $\mathcal I$ which contains the nonzero divisor $h$, makes that multiplier integral. Thus
\[
 \mathcal O_X\subset\mathcal F\subset\widetilde{\mathcal O}_X.
\]
We claim the quotient $\mathcal F/\mathcal O_X$ has exactly the nonnormal support. Only the reverse containment requires proof. At a nonnormal stalk let $k$ be the least integer with $\mathcal I_x^k\widetilde{\mathcal O}_{X,x}\subset\mathcal O_{X,x}$. It exists: the analytic Nullstellensatz and finite generation give $\mathcal I_x^\ell\subset(h)$ for some $\ell$, and $h\widetilde{\mathcal O}_{X,x}\subset\mathcal O_{X,x}$. Since $k>0$, choose $w\in\mathcal I_x^{k-1}\widetilde{\mathcal O}_{X,x}$ not holomorphic. Then $w\mathcal I_x\subset\mathcal O_{X,x}$, and every such product vanishes on $\{h=0\}$: $w$ is locally bounded on the regular locus, while its holomorphic factor tends to zero at each point of that set. Density and continuity of the holomorphic product give zero at the singular points as well. Thus $w\mathcal I_x\subset\mathcal I_x$, so $w\in\mathcal F_x\setminus\mathcal O_{X,x}$. This proves the claim and analyticity of the nonnormal locus.

For the last assertion, first extend the function across the analytic subset in $X_{\mathrm{reg}}$ using \cref{lem:bounded-holomorphic-removal}. It becomes weakly holomorphic and therefore holomorphic by normality.
\end{proof}
\index{complex space!nonnormal locus}

\Needspace{18\baselineskip}
\begin{theorem}[Normalization]\label{thm:analytic-normalization}
Every reduced complex space $X$ has a normal complex space $\widetilde X$ and a finite bimeromorphic map $\nu:\widetilde X\to X$, an isomorphism over $X_{\mathrm{reg}}$. This normalization is unique up to a unique holomorphic isomorphism over $X$. Its fiber at $x$ has one point for each irreducible branch of $(X,x)$, and
\[
 \nu_*\mathcal O_{\widetilde X}=\widetilde{\mathcal O}_X.
\]
\end{theorem}
\begin{proof}
First suppose $(X,x)$ irreducible. Choose finite module generators $f_1,\ldots,f_m$ of $\widetilde{\mathcal O}_{X,x}$ over $\mathcal O_{X,x}$, and a denominator $h$ such that $hf_j=g_j$ with $g_j$ holomorphic. On a small representative each $f_j$ satisfies a monic equation with holomorphic coefficients. In a local embedding $X\subset\Omega\subset\mathbb C^N$ form the graph of $(f_1,\ldots,f_m)$ off $\{h=0\}$ and take its closure $Y$ in $\Omega\times\mathbb C^m$. This closure is analytic. Indeed it is the closure of the complement of $\{h=0\}$ in the analytic set
\[
 z\in X,\qquad h(z)w_j=g_j(z).
\]
In a finite local irreducible decomposition of this set, keep exactly the branches on which $h$ is not identically zero. Their complements of $\{h=0\}$ are dense by parametrization and the identity theorem, so their finite union gives precisely that closure locally.

The monic equations of the $f_j$ hold on $Y$. They bound $w$ uniformly over compact base sets and allow only finitely many possibilities in each fiber. Therefore the projection $\nu:Y\to X$ is proper and finite. Over $X_{\mathrm{reg}}$ the graph map is holomorphic and its closure has just that graph, because each $f_j$ is bounded and holomorphic there. It is an isomorphism there, with dense image in $Y$. At $x$ the limits in \cref{prop:weak-holomorphic-denominator} show that the fiber is the singleton $y=(x,f_1(x),\ldots,f_m(x))$.

The germ $Y$ is normal at $y$. For a weakly holomorphic germ $a$ there, its pullback to the graph over $X_{\mathrm{reg}}$ is weakly holomorphic near $x$; properness and the singleton fiber give the required local bound. Write it $\sum_j a_jf_j$ with $a_j\in\mathcal O_{X,x}$. Lift the $a_j$ to ambient holomorphic functions. Then $\sum_j a_j(z)w_j$ is holomorphic near $y$ and equals $a$ on the dense graph, hence on the regular locus. Thus every weak germ there is holomorphic. By \cref{lem:normal-locus-open}, shrink around $y$ to normal points, and by properness shrink around $x$ so that its entire inverse image lies there. This gives a local normalization. For a reducible germ take the disjoint union of the local normalizations of its branches; its fiber at $x$ has the asserted number of points.

To prove uniqueness, an isomorphism between two normalizations is already prescribed over $X_{\mathrm{reg}}$. Near a point of the first normalization, take disjoint bounded coordinate neighborhoods of the finitely many points in the other fiber. Properness confines the image to their union after shrinking the base. The complement of the exceptional analytic set in a sufficiently small irreducible source neighborhood is connected: its connected dense covering locus remains connected after removing a proper analytic subset, by \cref{lem:bounded-holomorphic-removal}. Thus the image lies in one of those coordinate neighborhoods. Its bounded coordinate functions extend by \cref{lem:normal-locus-open}, and the target equations still hold by density. This extends the map holomorphically. Extend its inverse in the same way; their compositions are identities by density. These unique local maps glue, both giving global existence from the local constructions and uniqueness. Finally holomorphic functions on the normalization are weakly holomorphic downstairs by finite properness and the regular-locus identification; conversely a weak function pulls back to a bounded holomorphic function off the exceptional set and extends by normality. This proves the sheaf identity.
\end{proof}
\index{normalization!complex space}

\Needspace{9\baselineskip}
\begin{definition}\label{def:analytic-unibranch}
A reduced complex space is \emph{unibranch} if each of its local analytic germs is irreducible. Equivalently its local ring is a domain whose integral closure in its fraction field is local: \cref{prop:weak-holomorphic-denominator} identifies that closure, and on an irreducible germ its limit at the base point is the residue map, with units exactly the elements having nonzero limit. Thus this definition agrees with the local-ring definition of unibranchness in the analytic category.
\end{definition}
\index{complex space!unibranch}

\Needspace{11\baselineskip}
\begin{corollary}\label{cor:unibranch-normalization-topology}
On a unibranch space the normalization is a homeomorphism. Every point has arbitrarily small neighborhoods $U$ such that $U\setminus A$ is connected for every nowhere dense analytic subset $A$ of $U$.
\end{corollary}
\begin{proof}
The normalization has singleton fibers and is a proper continuous bijection, hence a homeomorphism since it is closed. For the second assertion choose an irreducible parametrization neighborhood. Its connected dense smooth covering locus remains connected after deleting $A$ by bounded holomorphic removal on that manifold. It is dense in $U\setminus A$, so that complement is connected too.
\end{proof}

\begin{example}\label{ex:cusp-node-normalization}
The curve $\{zw=0\}\subset\mathbb C^2$ has two branches at the origin; its normalization is the disjoint union of its two coordinate lines, with a two-point fiber there. It is not unibranch. The cusp $\{w^2=z^3\}$ is normalized by $t\mapsto(t^2,t^3)$. Indeed on its punctured curve $t=w/z$, and $t^2=z$, so this is an integral weakly holomorphic coordinate; the image is the whole cusp and its inverse is continuous at zero. The map is a homeomorphism but its inverse is not holomorphic as an ambient restriction: a power series in $t^2,t^3$ has no linear term. Thus the cusp is unibranch and nonnormal, distinguishing the two conditions.
\end{example}
\index{normalization!cusp}\index{normalization!crossing branches}

\section{Analytic curves and finite maps}\label{sec:analytic-curves-finite-maps}

An analytic curve detects a holomorphic function that varies on a singular germ. Finite parametrization supplies such curves without requiring the ambient space to be smooth.

\begin{proposition}\label{prop:analytic-curve-parametrization}
An irreducible one-dimensional reduced analytic germ $(C,x)$ has a parametrization $\gamma:(\Delta,0)\to(C,x)$ that is one-to-one and biholomorphic away from the origin. Its projection in suitable finite-parametrization coordinates is $t\mapsto t^q$ for some positive integer $q$. A normal one-dimensional complex space is a complex manifold.
\end{proposition}
\begin{proof}
Apply \cref{thm:analytic-local-parametrization} with one base coordinate $z$. Shrink its disc so that the discriminant has no zero except possibly zero. The covering of the punctured disc is connected, since the covering locus in that theorem is connected, and has $q$ sheets. A loop once around zero permutes these sheets in a single cycle. To justify this description directly, subdivide a path into finitely many intervals lying in evenly covered discs and lift it successively. Subdivision of a homotopy into small rectangles shows that homotopic paths with fixed endpoints give the same endpoint. Every loop in the punctured disc has an integer winding number: a continuously lifted argument along the loop changes by an integer multiple of $2\pi$, and subtracting that change contracts the loop in the logarithmic half-plane. Consequently every sheet reachable from a given one is obtained by powers of its one-loop permutation. Connectedness makes this permutation a single $q$-cycle.

On the half-plane $\{\operatorname{Re}\zeta<\log r\}$ lift $z=\exp\zeta$ holomorphically, starting on one sheet. The lift is independent of its path by the homotopy argument. It is periodic with period $2\pi q\mathrm i$; translation by $2\pi\mathrm i$ moves to the next sheet. Thus evaluating this lift at $q\log t$ defines a single-valued holomorphic map for $0<|t|<r^{1/q}$. Its ambient coordinates are bounded by the finite-parametrization neighborhood. One-variable removability extends them across zero. The defining equations extend there too, and the singleton fiber over $z=0$ makes $\gamma(0)=x$. The cyclic description proves injectivity and surjectivity on the punctured representatives. Rescale the parameter disc.

If $C$ is normal, the inverse parameter on its regular punctured neighborhood is bounded and therefore weakly holomorphic. Normality extends it holomorphically to $x$. The two compositions are identities on dense punctured neighborhoods and hence everywhere, proving that $x$ has a smooth disc coordinate.
\end{proof}
\index{analytic curve!parametrization}

\begin{proposition}\label{prop:analytic-curve-selection}
Let $(X,x)$ be a positive-dimensional irreducible reduced analytic germ, and let $(B,x)$ be a proper analytic subgerm. There is a nonconstant holomorphic map $\gamma:(\Delta,0)\to(X,x)$ whose punctured image avoids $B$ and lies in $X_{\mathrm{reg}}$.
\end{proposition}
\begin{proof}
Enlarge $B$ by the singular locus, and choose a holomorphic germ $f$ that vanishes on this enlarged set but not identically on $X$. In a finite parametrization over $\mathbb C^d$, the product of the values of $f$ on the generic sheets is a nonzero holomorphic base germ. Indeed it is bounded off the discriminant and extends by \cref{lem:bounded-holomorphic-removal}; it is nonzero because $f$ is a nonzero element of the finite field extension, so multiplication by $f$ is invertible over the base fraction field. Choose a complex line through zero on which neither this product nor the discriminant vanishes identically. Such lines exist by taking a direction where the lowest nonzero homogeneous term of their product is nonzero.

The inverse image of this line is analytic and finite over it. It has no positive-dimensional component contained in the fiber over zero, since that fiber is a singleton. Its components through $x$ that meet the punctured line are one-dimensional: on their regular loci a map with finite fibers has injective differential wherever its rank is locally maximal, by the constant-rank theorem, and its target has dimension one. At least one such component exists, since the line outside the two zero sets has nonempty fibers approaching $x$. Parametrize that component by \cref{prop:analytic-curve-parametrization}. After shrinking, its punctured points project outside both zero sets, and thus avoid $B$ and the singular locus.
\end{proof}
\index{analytic curve!selection}

\Needspace{7\baselineskip}
\begin{corollary}\label{cor:analytic-holomorphic-maximum}
If a holomorphic function on a reduced analytic germ attains a local maximum of its modulus at its base point, it is constant on each local irreducible branch. A compact analytic subset of an open subset of $\mathbb C^N$ is finite.
\end{corollary}
\begin{proof}
If $f-f(x)$ were nonzero on a branch, select a curve avoiding its zero set by \cref{prop:analytic-curve-selection}. The function $f\circ\gamma$ contradicts the one-variable maximum principle.

For the second assertion, suppose the analytic set has positive-dimensional points. Their set is closed and analytic locally: in the finite local decomposition retain the positive-dimensional branches; a zero-dimensional branch is an isolated point and cannot be a limit of the others without lying on them. This gives a nonempty compact analytic subset all of whose points have a positive-dimensional branch. At a maximum of $|z|^2$ on it, select a nonconstant analytic curve. The function $|\gamma(t)|^2$ is subharmonic, with Laplacian $4\sum_j|\gamma_j'(t)|^2$, and has a local maximum at zero. It must be constant, forcing every $\gamma_j'$ to vanish, a contradiction. The original set is therefore discrete and compact, hence finite.
\end{proof}

\begin{lemma}\label{lem:thin-set-path-connectivity}
If $E$ is relatively closed in a real ball $B\subset\mathbb R^N$, $N\geq2$, and $\mathcal H^{N-1}(E)=0$, then $B\setminus E$ is path connected.
\end{lemma}
\begin{proof}
Fix $a,b\in B\setminus E$. Intermediate points $c$ for which the segment $[a,c]$ meets $E$ belong to
\[
 \{a+t(e-a):e\in E,\ t\geq1\}.
\]
This set has $N$-dimensional Lebesgue measure zero. Exhaust $E$ by its intersections with compact subballs and restrict $t$ to bounded intervals. On each resulting product the displayed map is Lipschitz. A cover of the compact portion of $E$ by balls of radii $r_j$ with arbitrarily small $\sum_j r_j^{N-1}$ gives a cover of its product with a fixed interval by at most $C/r_j$ balls of comparable radius for each $j$, so its $N$-dimensional measure is zero. The countable exhaustion proves the claim. The same holds for $b$. Choose $c$ in the ball outside both null sets. Convexity keeps $[a,c]\cup[c,b]$ in $B$, and the choice keeps it disjoint from $E$.
\end{proof}

\begin{lemma}\label{lem:finite-analytic-open-map}
Let $(X,x)$ be a reduced pure $d$-dimensional analytic germ, and let $F:(X,x)\to(\mathbb C^d,0)$ be holomorphic with isolated fiber $F^{-1}(0)=\{x\}$ as a germ. It has a representative proper and finite over a ball about zero, and its image contains that ball.
\end{lemma}
\begin{proof}
For $d=0$ the assertion concerns a singleton. Otherwise embed the germ in a Euclidean neighborhood and take a small closed ball whose intersection with the zero fiber is just $x$. The compact boundary has positive minimum of $|F|$. Choose a target ball smaller than this minimum and restrict the source to its inverse image inside the open source ball. The map is proper there. Its fibers are compact analytic subsets of the source Euclidean neighborhood and hence finite by \cref{cor:analytic-holomorphic-maximum}.

On the regular locus the maximal rank is $d$: a smaller maximal rank would, on its open constant-rank locus, give positive-dimensional fibers. The union $C$ of the singular locus and the critical locus is a proper analytic subset of the pure $d$-dimensional source. Its analyticity can be expressed in an embedding by adjoining the rows of $\mathrm d F$ to the rows of differentials of local equations for $X$ and using minors, together with the equations of the singular locus. Its dimension is at most $d-1$ by \cref{prop:analytic-dimension-drop}. The image $E=F(C)$ is closed by properness, and $\mathcal H^{2d-1}(E)=0$: stratify $C$ by \cref{cor:analytic-strata-hausdorff}, cover the smooth strata by compact coordinate pieces, and use that $F$ is Lipschitz on each such piece.

Over the target ball minus $E$, the map is a local biholomorphism at every point of its source and proper. Its image is consequently both open and closed there. This image is nonempty: a regular full-rank point gives an open target image, which cannot be contained in $E$. The target complement is connected by \cref{lem:thin-set-path-connectivity}. The image therefore contains the entire complement and, being closed in the ball, contains the ball itself.
\end{proof}

\Needspace{9\baselineskip}
\begin{proposition}\label{prop:analytic-hypersurface-dimension}
On a reduced pure $d$-dimensional complex space, the zero set of a holomorphic function that is not identically zero on any local branch is either empty or pure of dimension $d-1$.
\end{proposition}
\begin{proof}
Dimension zero is immediate. At a zero $x$ on a branch $X_0$, dimension drop bounds the dimension of $V(f)$ by $d-1$. Choose $d-1$ ambient linear forms $\ell_j$, all zero at $x$, cutting its positive-dimensional branches successively until the common zero set of $(f,\ell_1,\ldots,\ell_{d-1})$ is isolated at $x$. At each step avoid the finitely many proper linear subspaces of forms that vanish identically on a remaining branch; dimension drop suffices, and any unused forms may be chosen generically. Apply \cref{lem:finite-analytic-open-map} to this $d$-tuple on $X_0$. Its image contains a target ball, so its zero set $f=0$ maps onto the target hyperplane in that ball. If its dimension were less than $d-1$, its image would have zero $(2d-2)$-dimensional measure, by the stratification and Lipschitz argument just used. This is impossible.

Apply this argument at a regular point of any branch of $V(f)$, away from the other branches of that zero set. A sufficiently small neighborhood has only this branch of the zero set, while every source branch remains $d$-dimensional. Thus this branch has dimension $d-1$. Every branch of $V(f)$ has the asserted dimension.
\end{proof}
\index{analytic set!hypersurface dimension}

\Needspace{20\baselineskip}
\section{Blowups and resolution}\label{sec:analytic-blowups-resolution}

Blowing up replaces a common zero set by the projective directions of its generators. The local construction proves the identities of ideals needed later on strict transforms.

\begin{definition}\label{def:analytic-blowup}
Let $X$ be reduced and let $\mathfrak a\subset\mathcal O_X$ be a coherent ideal nonzero on every local irreducible branch. On a neighborhood with generators $f_0,\ldots,f_r$, the \emph{blowup} $\operatorname{Bl}_{\mathfrak a}X$ is the reduced closure of the graph
\[
 x\longmapsto[f_0(x):\cdots:f_r(x)]
 \quad\hbox{in }X\times\mathbb P^r,
 \qquad x\notin V(\mathfrak a).
\]
Here $\mathbb P^r$ is the space of nonzero vectors in $\mathbb C^{r+1}$ modulo scalar multiplication, with its charts $\xi_i\ne0$ and coordinates $\xi_j/\xi_i$. It is compact: every projective point has a representative on the compact unit sphere, and the quotient map is continuous. The projection to $X$ is denoted $\pi$.
\end{definition}
\index{blowup!analytic ideal}

\Needspace{11\baselineskip}
\begin{proposition}\label{prop:analytic-blowup-universal}
The blowup is an analytic space, independent of the chosen generators. Its projection is proper and bimeromorphic, and $\mathfrak a\mathcal O_{\operatorname{Bl}_{\mathfrak a}X}$ is an invertible ideal. If a holomorphic map $g:Y\to X$ from a reduced space makes $\mathfrak a\mathcal O_Y$ an invertible ideal locally generated by a nonzerodivisor, then $g$ has a unique holomorphic lift to the blowup.
\end{proposition}
\begin{proof}
In each projective chart take the analytic set given by the equations of $X$ and $f_i\xi_j=f_j\xi_i$. The graph closure is obtained by retaining its local irreducible branches whose complement of $V(\mathfrak a)$ is nonempty. Finite local decomposition and density of the complement on those branches prove analyticity, as in the graph construction in \cref{thm:analytic-normalization}. It is closed in $X\times\mathbb P^r$, so compactness of projective space proves properness. Over the complement of $V(\mathfrak a)$ it is the graph and hence an isomorphism; that graph is dense in the blowup.

On $\xi_i\ne0$, the relations read $f_j=f_i(\xi_j/\xi_i)$, so the pulled-back ideal is generated by $f_i$. This generator is nonzero on every local branch of this chart: its nonzero graph points are dense there. It is thus a nonzerodivisor in the reduced local ring. This proves invertibility.

Locally on $Y$ write $f_i\circ g=d b_i$, where $d$ is a nonzerodivisor generator of the pulled-back ideal. The $b_i$ generate the unit ideal, and hence do not vanish simultaneously. They define the holomorphic projective map $[b_0:\cdots:b_r]$. Off $d=0$ this is the prescribed graph lift; density shows that its whole image lies in the graph closure. Different choices of $d$ give the same projective point. The lifts glue, and uniqueness follows from equality on the dense complement of $d=0$.

For two sets of generators, their blowups both make the original ideal invertible. The universal property gives maps between them, and their two compositions are identities by uniqueness. These canonical identifications also glue the local constructions on overlaps.
\end{proof}

An \emph{effective Cartier divisor} on a reduced space is a locally principal ideal generated by a nonzerodivisor, together with its local equations modulo multiplication by units. Its ideal is written $\mathcal O_X(-D)$, and the inverse invertible sheaf is $\mathcal O_X(D)$. Products of the ideals correspond to addition of effective Cartier divisors. Pullback or restriction composes the local equations with the map, provided the resulting equations remain nonzerodivisors. On a manifold a \emph{simple normal crossings divisor} has, in local coordinates, support $\{z_1\cdots z_k=0\}$ with its individual components smooth. Its multiplicities may be positive integers; simplicity refers to the reduced components and their intersections.
\index{divisor!effective Cartier}\index{divisor!simple normal crossings}

For a proper bimeromorphic map between manifolds, an effective divisor is called \emph{exceptional} if each of its components has image of codimension at least two. Proper images of analytic sets will be proved analytic in \cref{thm:proper-analytic-image}.

\Needspace{16\baselineskip}
\begin{proposition}\label{prop:smooth-center-blowup}
For a closed smooth center $C$ of positive codimension in a complex manifold $M$, the blowup $\pi:M'\to M$ is smooth. Its inverse-image divisor $E$ satisfies $\mathcal I_C\mathcal O_{M'}=\mathcal O_{M'}(-E)$, and its fiber over a point of $C$ is $\mathbb P^{r-1}$, where $r$ is the codimension of $C$. If $C$ is a divisor, the blowup is the identity and $E=C$.

Suppose that $W\subset M$ is a closed smooth submanifold and $C\subsetneq W$ is smooth. The strict transform $W'$, the closure of $\pi^{-1}(W\setminus C)$, is $\operatorname{Bl}_C W$, and
\begin{equation}\label{eq:strict-transform-ideal}
 \mathcal I_W\mathcal O_{M'}=\mathcal O_{M'}(-E)\mathcal I_{W'}.
\end{equation}
If the center is disjoint from $W$, its strict transform maps isomorphically to $W$ and the pulled-back ideal is $\mathcal I_{W'}$.
\end{proposition}
\begin{proof}
Choose coordinates $(z,w)$ with $C=\{w_1=\cdots=w_r=0\}$. The chart with $\xi_i\ne0$ has coordinates $t=w_i$, $v_j=w_j/w_i$ for $j\ne i$, and the unchanged $z$ coordinates. Its map is $w_i=t$, $w_j=tv_j$. These charts are smooth and fill the graph closure. They give $E=\{t=0\}$ and the stated projective fiber. For $r=1$ the coordinates and map are unchanged.

For the strict-transform identity take coordinates with $W=\{w_1=\cdots=w_r=0\}$ and $C=\{w_1=\cdots=w_r=z_1=\cdots=z_k=0\}$, $k\geq1$. Only charts whose pivot is one of $z_1,\ldots,z_k$ meet $W'$; there $w_j=t v_j$ and $W'$ is given by all $v_j=0$. Thus the left side of \eqref{eq:strict-transform-ideal} is $(t)(v_1,\ldots,v_r)$. In a chart whose pivot is a normal coordinate $w_i$, the strict transform is absent and the pulled-back ideal is $(t)$, since one of its generators is $t$. This proves the identity in all charts and identifies the restricted blowup with the charts of $\operatorname{Bl}_C W$. A center disjoint from $W$ has no effect in a neighborhood of $W$.
\end{proof}
\index{blowup!smooth center}\index{strict transform!ideal identity}

The following existence statements are the resolution and principalization input permitted in this book. Their source is \cite[Theorems 1.1 and 1.3, pp.~611 and 614; Definitions 2.1--2.2, pp.~618--619]{BM08}. The sentence extending those results to analytic spaces is on p.~611. The algorithms proving them are not reproduced here.

\begin{theorem}[Resolution of singularities]\label{thm:analytic-resolution-input}
A reduced complex space $X$ has a proper bimeromorphic map $\pi:M\to X$ with $M$ smooth, obtained by blowups with smooth centers and an isomorphism over $X_{\mathrm{reg}}$. The divisor formed by the inverse images of its blowup centers has simple normal crossings. The sequence is finite when $X$ is compact, and finite over each relatively compact open subset in general.
\end{theorem}

\begin{theorem}[Principalization]\label{thm:marked-order-one-input}
Let $M$ be a complex manifold, let $\mathfrak a\subset\mathcal O_M$ be a coherent ideal nonzero on each component, and let $D$ be a simple normal crossings divisor, possibly empty. There is a sequence of blowups $\sigma_i:M_i\to M_{i-1}$ with smooth centers, finite over relatively compact subsets, with the following properties. Set $M_0=M$ and $\mathfrak a_0=\mathfrak a$. Each center is contained in $V(\mathfrak a_{i-1})$ and has normal crossings with the accumulated divisor. If $E_i$ is its inverse-image divisor, then
\begin{equation}\label{eq:marked-order-one-transform}
 \mathfrak a_i=\mathcal O_{M_i}(E_i)\,
                   \mathfrak a_{i-1}\mathcal O_{M_i}
\end{equation}
is a holomorphic ideal, and the final ideal is the unit ideal. The accumulated divisor consists of the strict transforms of $D$ and all the $E_i$, and has simple normal crossings. The composite is an isomorphism off $V(\mathfrak a)$.
\end{theorem}

The second statement is the marked-ideal theorem with mark $1$ and $N=M$. Its cosupport is precisely $V(\mathfrak a_i)$, and its transformation divides by one exceptional factor, even when the ideal has larger vanishing order. On a compact input, or over a relatively compact base neighborhood where the sequence is finite, iterating \eqref{eq:marked-order-one-transform} gives
\[
 \mathfrak a\mathcal O_{M_t}=\mathcal O_{M_t}(-H),
 \qquad H=\sum_{i=1}^t(\sigma_{i+1}\circ\cdots\circ\sigma_t)^*E_i.
\]
Here the empty pullback at $i=t$ is $E_t$. The divisor $H$ is effective and has normal crossings support. For a coherent ideal on a reduced space nonzero on every local branch, apply resolution and then principalize its pullback, retaining the resolution boundary. This gives a \emph{log resolution}: on the smooth model the ideal is the ideal of an effective divisor, whose support together with the boundary divisor has simple normal crossings.
\index{resolution of singularities}\index{principalization!mark one}\index{log resolution}

\section{Proper images and relative principalization}\label{sec:proper-images-relative-principalization}

Properness prevents analytic branches from escaping at the boundary. The extension and mapping theorems below make the images of exceptional sets analytic; this property is needed when several blowups are composed.

\begin{theorem}[Remmert--Stein]\label{thm:remmert-stein-extension}
Let $p\geq0$ be an integer and $A$ an analytic subset of a reduced complex space $X$, of dimension at most $p$, and let $Z$ be an analytic subset of $X\setminus A$, all of whose local branches have dimension greater than $p$. Then its closure in $X$ is analytic.
\end{theorem}

\begin{theorem}[Remmert]\label{thm:proper-analytic-image}
The image of a proper holomorphic map of reduced complex spaces is analytic.
\end{theorem}
\begin{proof}
We prove the two statements together by induction. Write $R_m$ for the extension assertion with $\dim Z\leq m$, and $P_m$ for the proper-image assertion with source dimension at most $m$. Both assertions are local on the target, and a proper map has compact inverse image over a compact target neighborhood. Thus source dimensions and the number of source components meeting that inverse image are bounded there. The assertion $P_0$ follows since the image of a proper discrete space is closed and locally finite. The assertion $R_0$ has no nonempty case.

\emph{The extension step: $P_{m-1}$ implies $R_m$.} Embed a target neighborhood in $\mathbb C^N$. Induction on $p$ reduces to a smooth $p$-dimensional exceptional set: first extend across its regular locus, then across its singular locus, whose dimension is smaller by \cref{cor:analytic-strata-hausdorff}. For a nonpure $A$, use the union of its lower-dimensional branches and the singular loci and intersections of its top-dimensional branches as the smaller exceptional set. In suitable coordinates the smooth set is a linear $p$-plane $L$ through zero. Separating the finitely many possible dimensions of $Z$ reduces to pure dimension $s$, where $p<s\leq m$.

Choose linear forms $\xi_1,\ldots,\xi_s$ so that the successive intersections with $Z$ have pure dimensions $s-j$, and the first $p$ forms restrict independently to $L$. This is possible by \cref{prop:analytic-hypersurface-dimension}: at each stage avoid the forms identically zero on a branch of the preceding intersection. There are only countably many such branches, since the space has a countable basis and finite local decompositions. A countable union of proper linear subspaces has Lebesgue measure zero. We may at the same time arrange that $(\xi_1,\ldots,\xi_s)$ has full differential rank at a selected regular point of each component of $Z$: at each choice require independence of the restrictions of the forms at these countably many points. Subtract combinations of the first $p$ forms from the remaining forms to make them vanish on $L$. The joint zero intersections and the full-rank condition are unchanged by this triangular change. Complete them to ambient coordinates. Write $z=(z',z'')$, with $z'\in\mathbb C^s$, so that $L=\mathbb C^p\times\{0\}$, and let $\pi(z)=z'$.

The fiber over zero of $Z$ is discrete outside zero. Choose a small fiber ball $B''$ whose boundary misses this discrete set. The closed set $Z\cup A$ is disjoint from $\{0\}\times\partial B''$. Compactness of that boundary gives a base ball $B'$ with $(Z\cup A)\cap(B'\times\partial B'')=\varnothing$. Its projection to $B'$ is proper. Set $A'=B'\cap(\mathbb C^p\times\{0\})$. The restriction of $\pi$ from $Z$ over $B'\setminus A'$ is proper and finite, since its fibers are compact analytic subsets of $B''$.

Let $S_1$ be its singular source locus. Its proper image $S_1'$ is analytic in $B'\setminus A'$ by $P_{m-1}$. On the source over the complement of $S_1'$ let $S_2$ be the critical locus; its image $S_2'$ is analytic in $B'\setminus(A'\cup S_1')$ by the same induction. These images have dimensions at most $s-1$: their sources have those dimensions, and the stratification into smooth pieces with a Lipschitz map excludes a larger-dimensional analytic image. The complements are nonempty. The full-rank choice on every source component also ensures that $Z\setminus\pi^{-1}(A')$ is dense in $Z$.

Repeated bounded holomorphic removal shows that
\[
 G=B'\setminus(A'\cup S_1'\cup S_2')
\]
is connected and that every bounded holomorphic function on $G$ extends to $B'$. For connectedness, on a connected complex manifold the characteristic function of a disconnected complement of a proper analytic set would be a bounded holomorphic function; its extension would still take only the values zero and one and would be constant. Apply this first to $A'$, then to $S_1'$, and then to $S_2'$. Removal in the opposite order extends bounded functions. The finite projection over $G$ is an unramified covering onto $G$: its image is open and closed there by properness, and is nonempty by the full-rank condition and the dimension bounds for the exceptional images. Its sheet number is a fixed positive integer $q$.

For a linear form $\lambda$ in $z''$, define
\[
 P_\lambda(z',T)=\prod_{z\in Z\cap\pi^{-1}(z')}
                         (T-\lambda(z'')),\qquad z'\in G.
\]
The coefficients are bounded holomorphic functions and extend to $B'$. Each coefficient is a polynomial in the coefficients of $\lambda$ of its indicated degree; interpolation at finitely many fixed choices of $\lambda$ shows that its polynomial coefficients extend holomorphically as well. Thus the equations $P_\lambda(z',\lambda(z''))=0$ for all $\lambda$ are equivalent to finitely many coefficient equations. They define an analytic set $W$ in $B'\times B''$. On $G$, this set is precisely $Z$: a point outside a finite fiber can be separated from all its points by a linear form. Every fiber of $W$ is finite, even over the exceptional base. Indeed the coordinate forms bound each coordinate to finitely many roots, and any $q+1$ distinct points in one fiber could be separated by a linear form, contradicting the degree $q$ of its polynomial.

Near zero shrink to a proper representative of this finite projection, using a fiber ball whose boundary misses its finite zero fiber. Every local component of $W$ has dimension at most $s$, by the constant-rank theorem on its regular locus. A component of dimension $s$ maps onto a base neighborhood by \cref{lem:finite-analytic-open-map}; the same applies at each of its points. Its points over $G$ are consequently dense, since $G$ is dense in the base. A component of smaller dimension cannot meet $G$, where $W$ is a union of smooth $s$-dimensional covering graphs. Hence the closure of $W$ over $G$ is exactly the union of its $s$-dimensional components, and is analytic. Finally $Z$ over $G$ is dense in $Z$: the preimage of $A'$ is proper on each component by the full-rank choice, and over its complement the preimages of $S_1'$ and $S_2'$ are proper analytic subsets, since each component has a full-rank open image. Therefore this analytic union is precisely $\overline Z$. This proves $R_m$.

\emph{The mapping step: $R_m$ and $P_{m-1}$ imply $P_m$.} We may replace the source by its resolution from \cref{thm:analytic-resolution-input}; properness and surjectivity preserve its image and its dimension. This use is exactly the permitted resolution input. Work on each of the finitely many connected smooth components relevant over a compact target neighborhood, and embed that target in a Euclidean open set. Let $r$ be the maximal differential rank of the resulting map $F$ on one component. If $r=0$, its coordinate functions have zero differential and are constant on the connected component, so its image is a point. Otherwise the critical locus $C=\{\operatorname{rank}\mathrm d F<r\}$ is a proper analytic subset of source dimension at most $m-1$. Its image $F(C)$ is analytic by $P_{m-1}$ and has dimension at most $r-1$. For the latter assertion, stratify $C$ smoothly and then successively stratify the rank-drop loci of the restrictions of $F$. The constant-rank theorem covers the images of the resulting countably many pieces by submanifolds of dimension at most $r-1$, since the ambient differential has rank at most $r-1$ on $C$. An analytic image of larger dimension has a regular open part of positive $2r$-dimensional measure and cannot be such a union.

If $F(C)=F(X)$ the image is already analytic. Otherwise $F^{-1}(F(C))$ is a proper analytic subset of the connected manifold. Off it the rank is constantly $r$. A proper constant-rank map has locally analytic image: cover a compact fiber by finitely many constant-rank neighborhoods; their image germs are submanifolds, and properness confines the whole inverse image of a smaller target neighborhood to that finite cover. The image there is their finite union. Thus $F(X)\setminus F(C)$ is analytic and pure $r$-dimensional in the complement of $F(C)$. The theorem $R_m$ applies, since $\dim F(C)<r\leq m$. Its closure is the whole image: the source complement is dense, and the map is closed by properness. This proves $P_m$ and completes both inductions.
\end{proof}
\index{analytic set!Remmert--Stein extension}\index{proper map!analytic image}

\Needspace{7\baselineskip}
\begin{definition}\label{def:adapted-blowup-modification}
A composition of blowups of smooth centers on a complex manifold is \emph{adapted} to a closed smooth submanifold $W$ if at each step the center is either contained properly in the current strict transform of $W$ or disjoint from it. The identity composition is allowed.
\end{definition}
\index{modification!adapted to a submanifold}

\Needspace{21\baselineskip}
\begin{theorem}\label{thm:relative-principalization}
Let $M$ be a compact connected complex manifold, $W\subsetneq M$ a closed connected smooth submanifold, and $\mathfrak a\subset\mathcal O_M$ a coherent ideal whose restriction $\mathfrak a\mathcal O_W$ is nonzero. There is an adapted composition $f:M'\to M$, with strict transform $W'$, and effective divisors $H,H'$ on $M'$ such that $H$ is exceptional for $f$, $H'$ is disjoint from $W'$, and
\begin{align}
 \mathcal I_W\mathcal O_{M'}&=\mathcal O_{M'}(-H)\mathcal I_{W'},
 \label{eq:relative-center-ideal}\\
 \mathfrak a\mathcal O_{M'}&=\mathcal O_{M'}(-H-H'),
 \label{eq:relative-ideal-principalization}\\
 (\mathfrak a\mathcal O_W)\mathcal O_{W'}
     &=\mathcal O_{W'}(-H|_{W'}).
 \label{eq:relative-restricted-ideal}
\end{align}
The restriction $f|_{W'}:W'\to W$ is a composition of smooth-center blowups and an isomorphism outside a proper analytic subset of $W$.
\end{theorem}
\begin{proof}
Resolve the marked ideal $(\mathfrak b,1)$ on $W$, where $\mathfrak b=\mathfrak a\mathcal O_W$, by \cref{thm:marked-order-one-input}. Denote its successive smooth centers by $C_i$, its spaces by $W_i$, and its controlled ideals by $\mathfrak b_i$. Perform the same center blowups on the ambient spaces $M_i$. These centers are closed in $M_i$, since $W_i$ is closed; their ambient codimension is at least two, since they are proper in $W_i\subsetneq M_i$. The restricted maps and strict transforms are exactly the prescribed $W_i$ by \cref{prop:smooth-center-blowup}.

Start with $H_0=0$ and $\mathfrak a_0=\mathfrak a$, and set
\[
 H_i=\sigma_i^*H_{i-1}+E_i,\qquad
 \mathfrak a_i=\mathcal O_{M_i}(E_i)
                         \mathfrak a_{i-1}\mathcal O_{M_i}.
\]
These are holomorphic ideals and $\mathfrak a_i\mathcal O_{W_i}=\mathfrak b_i$. To check both assertions inductively, every generator of $\mathfrak a_{i-1}$ vanishes on $C_i$, because its restriction to $W_{i-1}$ belongs to $\mathfrak b_{i-1}$ and the center lies in the latter's zero set. Thus it belongs to $\mathcal I_{C_i}$ and its pullback is divisible by the single exceptional equation. The quotient remains holomorphic. In the tangential pivot charts meeting $W_i$, restriction of this quotient is exactly division by the exceptional equation for the blowup of $W_{i-1}$, including when the center is a divisor in $W_{i-1}$. This gives the controlled ideal with mark one, not division by its maximal order.

Iterating \eqref{eq:strict-transform-ideal} now gives
\[
 \mathcal I_W\mathcal O_{M_t}=\mathcal O_{M_t}(-H_t)\mathcal I_{W_t},
 \qquad
 \mathfrak a\mathcal O_{M_t}=\mathcal O_{M_t}(-H_t)\mathfrak a_t.
\]
At the last step $\mathfrak b_t=\mathcal O_{W_t}$. Hence $\mathfrak a_t$ is the unit ideal at each point of $W_t$: a generator nonzero after restriction is a unit in the ambient local ring. Its zero set is disjoint from $W_t$.

Principalize $\mathfrak a_t$ on $M_t$ using \cref{thm:marked-order-one-input} again, starting with empty boundary if necessary. All its centers lie over its zero set and are disjoint from the successive copies of $W_t$. Write this second composite as $g:M'\to M_t$, its total ideal divisor as $H'$, and set $H=g^*H_t$. Pulling back the last two identities proves \eqref{eq:relative-center-ideal} and \eqref{eq:relative-ideal-principalization}, and $H'$ is disjoint from $W'$. Restriction proves \eqref{eq:relative-restricted-ideal}. Every component of $H_t$ is exceptional in the ambient first sequence, whose centers have codimension at least two; its pullback remains supported on exceptional divisors of the whole map. Thus $H$ is exceptional. No such divisor contains $W'$, since the tangential charts show that the strict transform is not contained in $E_i$.

Finally the images in $W$ of all nontrivial centers of its restricted sequence are analytic by \cref{thm:proper-analytic-image} and have dimension smaller than $\dim W$. Their finite union is proper and contains the set outside which the sequence changes $W$. If $\dim W=0$, its nonzero ideal is the unit ideal and this first sequence is empty; the same reasoning applies with $H=0$.
\end{proof}
\index{principalization!relative}

\section{Connected fibers and resolution models}\label{sec:connected-fibers-models}

A proper bimeromorphic map can have positive-dimensional fibers. Unibranchness controls their connectedness, independently of smoothness or normality of the target.

\begin{theorem}\label{thm:unibranch-connected-fibers}
Let $f:Y\to X$ be proper and bimeromorphic between reduced complex spaces. If $X$ is unibranch, every fiber of $f$ is connected.
\end{theorem}
\begin{proof}
Choose a nowhere dense analytic subset $A\subset X$ outside which $f$ is an isomorphism. Its inverse image is nowhere dense in $Y$, by bimeromorphicity. Suppose the compact fiber over $x$ is disconnected, and partition it into two nonempty disjoint compact sets. Choose disjoint open neighborhoods $V_1,V_2$ of those sets in $Y$. Such neighborhoods exist by local compactness and the Hausdorff property, using finite coordinate covers of the two compact sets. Properness gives a neighborhood $U$ of $x$ whose whole inverse image lies in $V_1\cup V_2$: otherwise a sequence outside this union over base points tending to $x$ would have a convergent subsequence in the fiber.

Shrink $U$ to one of the neighborhoods in \cref{cor:unibranch-normalization-topology}, so that $U\setminus A$ is connected. Each $V_i\cap f^{-1}(U)$ is nonempty and open, and meets the dense set $Y\setminus f^{-1}(A)$. Therefore $f^{-1}(U\setminus A)$ is partitioned into two nonempty disjoint open sets by $V_1,V_2$. Its isomorphism with $U\setminus A$ contradicts connectedness.
\end{proof}
\index{proper map!connected fibers}

\begin{corollary}\label{cor:finite-bimeromorphic-normal-target}
A finite bimeromorphic map to a normal reduced complex space is an isomorphism. A holomorphic map from a normal space to a reduced space, whose inverse image of the target singular locus is nowhere dense, factors uniquely through the target normalization.
\end{corollary}
\begin{proof}
For the first assertion the target is unibranch. Its fibers are finite and connected, hence singletons. The proper map is a homeomorphism. Near its unique inverse point, choose bounded coordinate functions for the source. Composing with the inverse gives bounded holomorphic functions off a proper analytic subset of the target. They extend holomorphically by \cref{lem:normal-locus-open}, and their defining equations hold by density. Thus the inverse is holomorphic.

For the second assertion, lift the map uniquely over the inverse image of the target regular locus, using the normalization isomorphism there. Near a source point, properness of the finite normalization confines this lift to a union of disjoint bounded coordinate neighborhoods of its finitely many points over the image point. A sufficiently small irreducible source neighborhood minus its exceptional analytic set is connected by \cref{cor:unibranch-normalization-topology}; thus its lift lies in one of those coordinate neighborhoods. Its bounded coordinate functions extend on the normal source, and the target equations extend by density. These extensions give local lifts everywhere. Uniqueness on the dense original domain makes them glue and proves the asserted uniqueness.
\end{proof}
\index{normalization!factorization}\index{proper map!finite bimeromorphic}

\Needspace{10\baselineskip}
\begin{proposition}\label{prop:smooth-bimeromorphic-divisor-locus}
A proper bimeromorphic map $f:M'\to M$ between complex manifolds is an isomorphism outside an analytic subset of $M$ of codimension at least two. In particular every irreducible hypersurface in $M$ has a unique strict transform in $M'$ that maps isomorphically over its generic points.
\end{proposition}
\begin{proof}
Work on connected components of common dimension $n$. The critical locus of $f$ is analytic, given in coordinates by its Jacobian determinant; its image $S$ is analytic by \cref{thm:proper-analytic-image}. Outside $S$ the map is proper and a local biholomorphism, so its fibers are discrete and finite. They are singletons by \cref{thm:unibranch-connected-fibers}. The local inverses agree, giving an isomorphism there.

Suppose that $S$ has a hypersurface component $D$. Work near a regular point of $D$, away from the other components. Its inverse image $E$ is a pure hypersurface by \cref{prop:analytic-hypersurface-dimension}, since a local equation for $D$ is nonzero on the bimeromorphic source. Decompose $E$ locally over a compact target neighborhood into its finitely many relevant components. Remove from $D$ the images of components not dominating $D$, and, for each dominating component, the images of its singular locus and of the critical locus of its restriction to $D$. These are proper analytic subsets of $D$. Indeed the maximal rank on a dominating component is $n-1$, by the constant-rank stratification and the proper-image theorem; its singular and rank-drop loci have images of dimension at most $n-2$. The rank-drop locus together with the singular locus is closed analytic, as expressed by the defining-ideal and map Jacobian minors. Thus at a remaining point $a\in D$ the fiber is discrete in $E$ and hence finite.

Finite fibers persist over a small target neighborhood. To see this, put disjoint small Euclidean source balls around their finitely many points, with boundaries avoiding the fiber. Properness confines the inverse image of a small base neighborhood to these balls. On each boundary $|f-a|$ has a positive minimum; choose the base smaller than all those minima. The restricted maps are proper inside the balls, and all their fibers are compact analytic subsets of Euclidean neighborhoods, hence finite. The original map is consequently finite over this base neighborhood. It is an isomorphism there by \cref{cor:finite-bimeromorphic-normal-target}. This contradicts $a\in S$. Therefore $S$ has codimension at least two. Its complement meets every hypersurface densely; the closure of the inverse image of that dense part defines its unique strict transform. This closure is analytic by retaining the branches of the analytic inverse-image hypersurface that are not contained over $S$.
\end{proof}

A \emph{smooth model} of $X$ is a proper bimeromorphic map $M\to X$ with $M$ smooth. Resolution gives such models. Two models have a common smooth model: in $M_1\times_X M_2$ take the closure of their common graph over an open set where both maps are isomorphisms, and resolve that closure. The fiber product is analytic, since equality of the two images is expressed by coordinate equations in local embeddings of $X$. The graph closure is analytic by retaining the local components not contained over the exceptional analytic set, as in \cref{prop:analytic-blowup-universal}. Its projections are proper: over a compact subset of $M_1$, the second coordinates lie in the compact inverse image under $M_2\to X$ of its image in $X$, and the fiber product is closed in that compact product. Their images are dense over the common open set and closed, hence surjective. The resolved projections are bimeromorphic.

When $X$ is compact, there is also a common refinement by smooth-center blowups when the specified models themselves are finite smooth-center blowup sequences over $X$. Begin with a resolution of $X$. To lift the current smooth model through the next blowup in either specified sequence, principalize the pulled-back center ideal using \cref{thm:marked-order-one-input}, and use \cref{prop:analytic-blowup-universal}. Repeat for the finitely many centers, and then for the second sequence. Each pulled-back center ideal is nonzero on each relevant component, since it is the unit ideal over the common isomorphism locus. The result is still a finite smooth-center sequence over $X$ and dominates both specified models. This refinement statement concerns the specified blowup sequences; it uses no assertion that every proper bimeromorphic map is itself such a sequence.
\index{resolution!common refinement}

\begin{definition}\label{def:prime-divisor-model}
A \emph{prime divisor over $X$} is a pair consisting of a smooth model $M\to X$ and an irreducible hypersurface $F\subset M$, where pairs are identified if their strict transforms agree on a common smooth model. By \cref{prop:smooth-bimeromorphic-divisor-locus} strict transforms preserve the generic hypersurface, and common refinements make this an equivalence relation. The divisor is \emph{exceptional} if its image in $X$ has codimension at least two. One may choose a model on which its representative is smooth: principalize its ideal and retain its strict transform, a component of the resulting simple normal crossings support.
\end{definition}
\index{prime divisor!over a complex space}

\begin{theorem}\label{thm:divisor-smooth-blowup-realization}
Let $X$ be an irreducible compact reduced complex space. Every prime divisor over $X$ is represented on a smooth model obtained from $X$ by a finite sequence of blowups with smooth centers. Its representative may also be chosen smooth.
\end{theorem}
\begin{proof}
First resolve $X$ by \cref{thm:analytic-resolution-input}. Taking a common smooth model with a model representing the given divisor reduces the assertion to a compact complex manifold $X$. Indeed strict transforms preserve the given divisor by \cref{prop:smooth-bimeromorphic-divisor-locus}, and a finite blowup sequence over this resolved $X$ can be composed with the initial resolution. We therefore have a proper bimeromorphic map $f:M\to X$ of compact manifolds and a prime divisor $E\subset M$.

The local Jacobian determinant $J_f$ is holomorphic and nonzero on the isomorphism locus. Its vanishing order along $E$ is a finite nonnegative integer $k$. Changes of coordinates multiply $J_f$ by holomorphic units, so this order is well defined. For composable maps the chain rule gives
\begin{equation}\label{eq:divisor-jacobian-order-chain}
 \operatorname{ord}_E J_{g\circ h}
 =\operatorname{ord}_E J_h+
              \operatorname{ord}_E(J_g\circ h).
\end{equation}
Replacing $M$ by a common smooth model does not change $k$: the new map to $M$ is an isomorphism at generic points of the strict transform of $E$, again by \cref{prop:smooth-bimeromorphic-divisor-locus}.

The center $Z=f(E)$ is a proper irreducible analytic subset by properness and \cref{thm:proper-analytic-image}. If it is a hypersurface, $f$ is an isomorphism over a dense open subset of $Z$, and $Z$ already represents $E$. Suppose instead that $Z$ has codimension at least two. Its reduced ideal is coherent by \cref{thm:analytic-ideal-coherence}. Principalize that ideal on $X$ by \cref{thm:marked-order-one-input}, obtaining a finite smooth-center sequence $\sigma:X'\to X$ such that $\mathcal I_Z\mathcal O_{X'}$ is invertible. Take a common smooth model of $M$ and $X'$; write its maps as $h:M'\to M$ and $f':M'\to X'$, and let $E'$ be the strict transform of $E$.

At least one nontrivial blowup center in this sequence contains the current generic center of $E'$. To check this, follow the sequence near a generic point of $Z$ until the first such center. Earlier centers avoid a dense open subset of $Z$, so their blowups are isomorphisms there and the current ideal is still the reduced ideal of a set of codimension at least two. If there were no first center, that ideal would remain unchanged near a generic point of $Z$ and could not become invertible. A nonzero invertible proper ideal on a manifold defines a hypersurface, by \cref{prop:analytic-hypersurface-dimension}. Moreover the first center has codimension $r\geq2$: it is contained in the current zero set, which still agrees with $Z$ there. Centers with different components can be treated on their component containing this generic point.

In the smooth blowup chart of \cref{prop:smooth-center-blowup}, the normal coordinates are $w_i=s$ and $w_j=sv_j$ for $j\ne i$. The Jacobian determinant is a unit times $s^{r-1}$. Pulling its exceptional equation $s$ back to $M'$ gives positive order along $E'$, because the image of $E'$ lies over that center. It is not identically zero on $M'$, since the maps are bimeromorphic. All the other blowup Jacobians have nonnegative order. Thus
\[
 \operatorname{ord}_{E'}(J_\sigma\circ f')\geq1.
\]
Applying \eqref{eq:divisor-jacobian-order-chain} to $f\circ h=\sigma\circ f'$ yields
\[
 0\leq\operatorname{ord}_{E'}J_{f'}
       =k-\operatorname{ord}_{E'}(J_\sigma\circ f')\leq k-1.
\]
Consequently each repetition with a center of codimension at least two strictly decreases a nonnegative integer. After at most the initial $k$ repetitions, the center is a hypersurface on the current smooth model. This proves the required finite realization. Finally principalize its reduced ideal while retaining the accumulated boundary. The total transform has simple normal crossings support, so its strict transform is smooth and still represents the same divisor. In dimension zero there are no prime divisors; in dimension one the center is already a hypersurface.
\end{proof}
\index{prime divisor!realization by smooth blowups}

\section{Plurisubharmonic functions on reduced spaces}\label{sec:psh-reduced-spaces}

At a singular point there are fewer ambient directions inside the space. The definition retains ambient plurisubharmonicity; disc pullbacks are a consequence of that definition.

\begin{definition}\label{def:psh-reduced-space}
A function $u:X\to[-\infty,+\infty)$ on a reduced complex space is \emph{locally plurisubharmonic} if in a local embedding $X\cap U\subset\Omega\subset\mathbb C^N$ it is the restriction of a plurisubharmonic function on an ambient neighborhood, allowing the constant function $-\infty$. Write $u\in\operatorname{PSH}(X)$ when in addition it is not identically $-\infty$ on any irreducible component. Local restriction properties will also be used when a component is identically $-\infty$.
\end{definition}
\index{plurisubharmonic function!on a reduced space}

This definition is independent of the embedding. If two embeddings have coordinates $z$ and $w$, the functions $w_j$ on the space extend holomorphically to the $z$-ambient neighborhood by the definition of $\mathcal O_X$. Their joint extension maps a smaller neighborhood into the $w$-ambient domain. Pulling back an ambient psh extension by this holomorphic map gives the required extension for the $z$ embedding, by \cref{prop:psh-holomorphic-pullback}. The converse follows by interchanging the embeddings. The same argument proves holomorphic pullback on reduced spaces. In particular $u\circ\gamma$ is subharmonic or identically $-\infty$ for every holomorphic disc $\gamma$ in $X$.

\Needspace{9\baselineskip}
\begin{proposition}\label{prop:psh-compact-connected-constant}
A locally plurisubharmonic function on a compact connected reduced complex space is constant, possibly $-\infty$.
\end{proposition}
\begin{proof}
The function is upper semicontinuous and attains its maximum $M$. If $M=-\infty$, the assertion is immediate. At a point attaining a finite maximum, restrict to a small irreducible parametrization neighborhood on each local branch. A curve from \cref{prop:analytic-curve-selection} through this point and meeting the smooth covering locus gives a subharmonic disc function with a maximum at zero. It is constant, so a covering-locus point also attains $M$. The usual psh maximum principle then makes $u$ equal to $M$ on the connected covering locus. This principle follows from the subharmonic maximum principle on a manifold, or successively from its coordinate-disc submean inequalities. For any other point of the branch, choose a curve through it whose punctured image avoids the discriminant preimage and the singular locus. Its punctured values are $M$, so the subharmonic mean inequality and upper semicontinuity give its value $M$ at the center. On a zero-dimensional branch the sole value already is $M$. Thus the maximum set is open. It is closed, since $\{u<M\}$ is open for an upper semicontinuous function bounded by $M$. Connectedness gives the conclusion.
\end{proof}

\begin{proposition}\label{prop:singular-codimension-two-upper-bound}
Let $A$ be an analytic subset of a reduced complex space $X$ with codimension at least two on every local branch. Every locally plurisubharmonic function on $X\setminus A$ is locally bounded above near $A$.
\end{proposition}
\begin{proof}
It suffices to work on one irreducible parametrization neighborhood at a point of $A$. Its dimension $d$ is at least two. Let $\pi:X\to B\subset\mathbb C^d$ be its finite proper projection with singleton central fiber. The image $\pi(A)$ is analytic by \cref{thm:proper-analytic-image} and has dimension at most $d-2$ by the stratification and Lipschitz estimate. Choose a base line through zero whose intersection with this image is isolated at zero; avoid the nonzero lowest homogeneous term of an equation of the image, and then shrink the line disc. Write base coordinates $(\zeta,z')$ in its direction. For some $r>0$ and $\varepsilon>0$, the cylinder
\[
 \{|\zeta|=r,\ |z'|\leq\varepsilon\}
\]
is compactly contained in $B$ and disjoint from $\pi(A)$. Its inverse image is compact and disjoint from $A$, so the given function has an upper bound $M$ there.

Projection of $\pi(A)$ from $\{|\zeta|<r,\ |z'|<\varepsilon\}$ to the $z'$ ball is proper: its closure misses the fiber-circle boundary. Its image $S$ is analytic and has dimension at most $d-2$. For $z'\notin S$, the curve $C_{z'}=\pi^{-1}(\{|\zeta|<r\}\times\{z'\})$ avoids $A$. It is pure one-dimensional. Indeed, cut successively by the fixed $d-1$ base coordinates. At each stage a local branch maps finitely and openly to the remaining base coordinates, by \cref{lem:finite-analytic-open-map}; the next coordinate equation is therefore nonzero on that branch. The dimension drops by one at each cut by \cref{prop:analytic-hypersurface-dimension}. Its normalization is a smooth curve by \cref{prop:analytic-curve-parametrization}, finite and proper over the base disc. Pullback of $u$ to that normalization is subharmonic, allowing $-\infty$. It is at most $M$ by the maximum principle: its boundary over $|\zeta|=r$ is bounded by $M$, and a larger interior maximum would make it constant on a component. Each such component reaches the boundary, since a finite proper map of a curve to a disc is open by \cref{lem:finite-analytic-open-map}. To apply this argument literally on a compact set, take the normalization over a slightly larger line disc; the compact inverse image of $|\zeta|\leq r$ then lies inside its domain. This inverse image is disjoint from $A$, so the pulled-back function is defined on a neighborhood of it.

This bounds $u$ on the complement of $\pi^{-1}(S)$ in the cylinder. At any point $y\notin A$ over $S$, choose on each local branch a curve through $y$ avoiding $A\cup\pi^{-1}(S)$ off its center. The set $\pi^{-1}(S)$ is proper on every local branch: finite projections on a pure $d$-dimensional branch are open by \cref{lem:finite-analytic-open-map}. The curve selection theorem therefore applies. On its punctured disc $u\leq M$, and the submean inequality gives $u(y)\leq M$ too. The finitely many original branches give a local upper bound on $X\setminus A$ near the chosen point.
\end{proof}
\index{plurisubharmonic function!codimension-two upper bound}
